%% This is Fortran lecture notes, for SMA HPCES
%% complete self-contained (no include, external macros)
% TeX Macro for lecture notes (modified 24 Sep 1999)
% dimensions
\hsize 159.2 mm
\vsize 242.2 mm
\parindent 0pt
\parskip 4pt plus 2pt minus 1pt
\overfullrule 0pt
% fonts
\font\titlefont=cmbx10 scaled 1440
\font\lecfont=cmbx12
\font\sc=cmcsc10
\font\ninerm=cmr9
\font\ninett=cmtt9
\def\h{\hskip 1cm}
\def\t{\hskip 6mm}
% Define the verbatim program listing macro \listing{filename}
%
\def\uncatcodespecials{\def\do##1{\catcode`##1=12 }\dospecials}
\def\listing#1{\par\begingroup\setupverbatim\input#1 \endgroup}
\def\setupverbatim{\tt 
  \def\par{\leavevmode\endgraf} \catcode`\`=\active
  \obeylines \uncatcodespecials \obeyspaces
  \everypar={\t}
  \parskip=-1pt plus 1pt minus 1pt
}
{\obeyspaces\global\let =\ } % let active space = \control space
{\catcode`\`=\active \gdef`{\relax\lq}}
%
% verbatim text macro with smallskip before and after
% \ttbg
% verbatim text
% \tted
\def\ttbg{\par\smallskip\begingroup\setupverbatim\doverbatim}
{\catcode`\|=0 \catcode`\\=12 % | is temporary escape character
  |obeylines|gdef|doverbatim^^M#1\tted{#1|endgroup|smallskip}}
%
% one line verbatim
% \v`text`
\def\v{\begingroup\t\tt\uncatcodespecials
  \obeyspaces\doverbatt}
\newcount\balance
{\catcode`<=1 \catcode`>=2 \catcode`\{=12 \catcode`\}=12
  \gdef\doverbatt`<\balance=1\verbatimloop>
  \gdef\verbatimloop#1<\def\next<#1\verbatimloop>%
    \if#1`\advance\balance by-1
     \ifnum\balance=0\let\next=\endgroup\fi\fi\next>>

% the one line verbatim with comment in 9pt rm
%\l{\v`test`\m{comment}}
\def\m#1{\hfill {\ninerm #1} \t}
\def\l{\hbox to \hsize}
% one line of tt font without comment
\def\ll#1{\leftline{\t\tt #1}}
\def\ttsp{\hphantom{2}}
\newskip\hugeskipamount
\hugeskipamount = 20mm plus3mm minus3mm
\def\hugeskip{\vskip\hugeskipamount}
\def\hugebreak{\par \ifdim\lastskip<\hugeskipamount
  \removelastskip \penalty-500 \hugeskip \fi}
\def\lecture#1{\hugebreak {\noindent\lecfont Lecture #1} \bigskip\nobreak}
\def\section#1{\medbreak {\noindent\sc #1} \smallskip\nobreak}

\def\{{\char`\{}
\def\}{\char`\}}
\def\sms{\smallskip}
\def\fninety{$^{\scriptscriptstyle 90}$}

%def \problem and \answer with automatic numbering
\newcount\probno
\probno=0

\outer\def\problem{\medbreak\smallskip \global\advance\probno by 1
\item{{\bf\the\probno}.}}

\chardef\other=12
\newwrite\ans
\immediate\openout\ans=answers %file for answers to problems
\outer\def\answer{\par\medbreak
  \immediate\write\ans{}
  \immediate\write\ans{\string\ansno{\the\probno}}
  \copytoblankline}
\def\copytoblankline{\begingroup\setupcopy\copyans}
\def\setupcopy{\def\do##1{\catcode`##1=\other}\dospecials
    \catcode`\|=\other \obeylines}
{\obeylines \gdef\copyans#1
  {\def\next{#1}%
  \ifx\next\empty\let\next=\endgroup %
  \else\immediate\write\ans{\next} \let\next=\copyans\fi\next}}
\def\ansno#1{\medbreak\noindent{\bf #1:\ }}

%end of macro definitions

\def\ttsp{\hphantom{2}}
\def\]{\leavevmode\hbox{\tt\char`\ }}
{\nopagenumbers
\centerline{\titlefont FORTRAN PROGRAMMING}
\vskip 2cm
\centerline{Dr.~WANG Jian-Sheng}
\medskip
\centerline{Computational Science, Rm
S17--07--21, tel x6880, wangjs@cz3.nus.edu.sg}
\bigskip
\centerline{October 1999}
\vskip 2cm 
\leftline{References: {\sl Programmer's Guide to Fortran 90}, 
W.~S.~Brainerd, C.~H.~Goldberg, \& J.~C.~Adams}
\smallskip
\leftline{\hphantom{References:} {\sl Fortran 90 Explained},
M.~Metcalf \& J. Reid}
\leftline{\hphantom{References:} {\sl Using MPI}, W. Gropp, E. Lusk, \& 
A. Skjellum}
\bigskip
\vfill
\centerline{\titlefont Contents}
\bigskip

\line{Lecture \hfil Topics \hfil Chapter}
\smallskip
\l{\t\ttsp 1.\h History, Examples\hfill 1\t}
\l{\t\ttsp 2.\h Data Types, Expressions, and Assignments \hfill 1\t}
\l{\t\ttsp 3.\h Control Statements \hfill  2\t}
\l{\t\ttsp 4.\h Basic Program Units and Procedures \hfill 3\t}
\l{\t\ttsp 5.\h Input, Output, and Files \hfill 9\t}
\l{\t\ttsp 6.\h Arrays \hfill 4\t}
\l{\t\ttsp 7.\h Character Data \hfill 5\t}
\l{\t\ttsp 8.\h Pointers \hfill 8\t}
\l{\t\ttsp 9.\h Structures, Derived Types, and Linked Lists \hfill 6\t}
\l{\t 10.\h Recursion  and Trees\hfill 7\t}
\l{\t 11.\h Modules \hfill 6\t}
\l{\t 12.\h Intrinsic Procedures \hfill appendix B\t}
\l{\t 13.\h Introduction to MPI \hfill \t}

\vfill\eject
}
\lecture{1. History, Examples}

\section{History of Fortran}

In the early days of computer programming, the only way to write a program
is in assembly codes.   Fortran (from FORmula TRANslation) was
developed around 1955 at IBM to ease the tedious burden of using
assembly language.  It is the first high-\-level language.
Since then, Fortran has become the language of
choice for scientific computation.  Here is a typical reaction of
using Fortran in the early days:
\smallskip
{\parindent 1cm\narrower
\noindent The programmer attended a one-day course on Fortran and spent 
some more
time referring to the manual.  He then programmed the job in four hours,
using 47 Fortran statements.  These were compiled by the IBM 704 in
six minutes, producing about 1000 instructions.  He ran the program
and found the output incorrect.  He studied the output and was able
to localize his error in a Fortran statement he had written.  He rewrote
the offending statement, recompiled, and found that the resulting
program was correct.  He estimated that it might have taken three days
to code the job by hand, plus unknown time to debug it, and no appreciable
increase in speed of execution would have been achieved thereby.\par}

\smallskip
The language itself has evolved
and changed a lot ---  from Fortran 66 to Fortran 77, and to more recent
Fortran 90/95.   Recent new standard is a big jump from an old fashioned
to a new modern-looking language.   Here are few strong points to use
Fortran: (1) tradition for scientific computation; (2) many library
routines, e.g. NAG, IMSL for scientific computation; (3) most efficient
and optimized compilers on supercomputers (e.g., Cray); (4) new, powful
vector processing features in Fortran 90.   On the other hand, more
and more people are now using C or C++ for their good integrations with
UNIX operating system and other advantages.

\section{First example}
We shall study Fortran~90 most of the time.  Since there are still many
legacy Fortran 77 code around, and many of the older generations till
program in Fortran 77, we'll also mention Fortran~77.
Fortran~90 includes all of the Fortran~77 standard with many
new features.

Here is a very simple example of Fortran program that will compile and run.
\sms
\ll{\]\]\]\]\]\]}
\l{\v`C     print N Hellos`\m{comment starts with {\ninett C} at column 1}}
\l{\v`      INTEGER I, N`\m{declarations and statements start from column 7}}
\sms
\l{\v`      PRINT *, 'enter an integer'`\m{write out at screen}}
\l{\v`      READ *, N`\m{read from keyboard}}
\sms
\l{\v`      DO 100 I = 1, N`\m{a loop, like {\ninett for} in C}}
\l{\v`100      PRINT *, I, ' Hello'`\m{{\ninett 100} is 
				Fortran 77 style loop label}}
\sms
\l{\v`      END`\m{end program unit}}
\sms
The executable is obtained with
\sms
\ll{f90 $file$.f90}
\sms
for the Fortran 90 compiler or with
\sms
\ll{f77 $file$.f}
\sms
for Fortran 77 compiler.  Please note the standard extensions for
file names.


\section{Source form}

(1) Fixed Source Form


The above example is in Fortran~77 fixed source form -- 
The letter {\tt C} at column 1 starts a comment line;
columns 1 to 5 are for
labels; columns 7 to 72 are for statements; anything except space or zero
at column 6 indicates a continuation of the previous line.  Each
statement occupies one line.
\sms

(2) Free Source Form

In Fortran the line has significance.  This is also the case for Fortran~90.
Fortran~90 allows a free source form.  Specifically,
each line may contain up to 132 characters.
A statement can start anywhere and continue to the next line with
ampersand ({\tt\&}).   Semicolons separate multiple statements on a single 
line.  The exclamation mark ({\tt !}) starts a comment to the end of line.
E.g.,
\ttbg
p = n * R * T / V             !  comment: ideal gas law
y = a * x**2 +       & 	      !  continued to next line
    b * x + c
A = 0; B = 0; C = 0           ! multiple statements on a single line
LONG_STRING =       
   'String is quoted in apostrophes &
   &and continued in several lines like&
   & this.'
\tted

\section{Names (or identifiers)}
Fortran names are not case sensitive.  Uppercase and 
lowercase are treated as the same.  A name can be up to 31 characters from
letter {\tt A} to {\tt Z}, underscore {\tt \_}, and digit {\tt 0} to {\tt 9}.
The first character must be a letter.  Fortran 77 requires the length of the
names not exceeding six.  Here are some valid Fortran names:
\ttbg
MASS
ACT_1
REAL
X3J3
\tted

\section{Second example}
The following page contains a program (in Fortran 90 style) 
that calculate the roots of a 
quadratic equation with real coefficients.  But the roots can be complex.
Thus, the use of complex variables are especially convenient.
\ttbg
PROGRAM ROOTS
!  solve the equation   A * x**2 + B * x + C = 0
   IMPLICIT NONE
   REAL :: A, B, C
   COMPLEX :: X1, X2, D

   READ *, A, B, C
   PRINT *, "Input data A:", A
   PRINT *, "           B:", B
   PRINT *, "           C:", C

   D = B**2 - 4 * A * C
   X1 = ( - B + SQRT(D) )/(2 * A)
   X2 = ( - B - SQRT(D) )/(2 * A)

   PRINT *, "The roots are:"
   PRINT *, "x1 = ", X1
   PRINT *, "x2 = ", X2

END PROGRAM ROOTS
\tted

\vfill\eject
\lecture{2. Data Types, Expressions, and Assignments}

\section{Basic data types}
Variable declarations take the form $type$ [, $attribute$ (optional)] 
{\tt ::} $variable$, e.g.,
\sms
\l{\v`INTEGER  :: I, J, K`\m{the double colon is optional}}
\l{\v`REAL :: X`\m{floating point number}}
\l{\v`LOGICAL :: LVALUE = .TRUE.`\m{initialization in declaration}}
\l{\v`COMPLEX :: C1, C2 = (0.0, 1.0)`\m{complex number is represented by pair of reals}}
\l{\v`CHARACTER (LEN=20) :: NAME`\m{store 20 characters}}
\l{\v`INTEGER, DIMENSION(1:10) :: A`\m{declare an array}}
\l{\v`REAL, POINTER :: P`\m{a pointer to real}}
\l{\v`TYPE (STRUCT) :: TMP`\m{user-defined type}}
\sms
Fortran 77 does not allow initialization during declaration and
there should be not double colon.   The complex data type is convenient for 
scientific calculations.  The complex value is represented by a pair
of real numbers, enclosed in parentheses, like
\sms
\ll{(2.0, 5.0)}
\sms
which represents the complex number $ 2.0 + 5.0 i$.

Integer and real variables need not declared.  When a variable is not
declared, its type is implied through its name.  Variable names
starting with {\tt I}, {\tt J}, {\tt K}, {\tt L}, {\tt M}, or {\tt N}
are assumed of integer type.  Otherwise, it is assumed real. You can
change the rule by specifying 
\sms
\ll{IMPLICIT INTEGER (I-K, M, N), LOGICAL (L), COMPLEX(C)}
\sms
This means that variables with names starting with the letter {\tt I}, 
{\tt J}, {\tt K}, {\tt M}, {\tt N}, are integers,  variables
beginning with {\tt L} are logical, and complex if the first letter
is {\tt C}.  Anything else is real. 
Implicit type is handy for short programs, but it could be a source
of programming errors. To prohibit undeclared variables, one can put
\sms
\ll{IMPLICIT NONE}
\sms
at the beginning of each program unit.

\section{Kind parameter}
In Fortran~90, each variable or literal constant has a {\sl kind} 
as well as a type. 
It is possible to specify different kinds (size or precision) of real
or integer by a kind parameter, e.g.
\ttbg
INTEGER, PARAMETER :: LONG = SELECTED_REAL_KIND(9, 99)
REAL (KIND = LONG) :: A, B
\tted
The first line declares {\tt LONG} as an integer with {\tt PARAMETER}
attribute.  It means that {\tt LONG} is a constant and its value can
not be changed. 
The intrinsic function {\tt SELECTED\_REAL\_KIND} returns an integer (as
a constant, or parameter) which represents a type of real with
at least 9 significant decimals and
with range between $10^{-99}$ to $10^{99}$.  Whether the type of real
requested exists is another question.  If the required type is not available,
the function will return a negative number.  The function {\tt
SELECTED\_INT\_KIND($R$)} returns the kind parameter value for an
integer data type in the range $-10^R$ to $10^R$.    In the above
declaration, {\tt A} and {\tt B} are real of the kind with at least 9
digits precision and with a range at least from $10^{-99}$ to 
$10^{99}$.  One can specify
literal constants with a particular {\tt KIND}, e.g.,
\sms
\ll{36\_2\h 1.0\_LONG \h  1\_'first kind'}
\sms
Note for character string, the kind parameter is prefixed. 

Another intrinsic function associated with the kind parameter of data
type is {\tt KIND($x$)}.  {\tt KIND(1)} is the default integer kind;
{\tt KIND(0.0)} is the default real kind; and {\tt KIND('A')} is the
default character kind.

In Fortran~77, kind parameter is not available.  Instead, one can
usually specify the size of the data with type name like 
\sms
\ll{INTEGER*2 {\rm ,\quad} INTEGER*4 {\rm ,\ \ or \quad} REAL*8}
\sms
which represent two-byte, four-byte integers, eight-byte reals,
respectively. 

\goodbreak
\section{Arithmetic expressions}
The notations for arithmetic operations are standard except the 
exponentiation operator {\tt **}, e.g.
\sms
\l{\v`A * X**Y + B/Y - 2`\m{in math notation $a x^y + {b/y} -2$}}
\l{\v`A ** B ** C`\m{$A^{B^C}$}}
\l{\v`A / B / C`\m{$(A/B)/C$}}
\l{\v`8 / 3`\m{2, integer division truncates towards 0}}
\l{\v`2 ** 3`\m{8}}
\l{\v`2 ** (-3)`\m{result is 0}}
\sms
Assignments between different types are allowed, but cause 
conversion or promotion.  E.g., if {\tt I} is an integer and {\tt A} is
a real, then
\sms
\l{\v`I = 7.3`\m{{\ninett I} gets value {\ninett 7}}}
\l{\v`A = (4.019, 2.217)`\m{{\ninett A} gets real part {\ninett 4.019}}}
\sms
Conversion can be made explicitly with the functions
\sms
\ll{INT($x$)}
\ll{REAL($x$)\t {\rm or}\t REAL($x$, $kind$)}
\ll{CMPLX($x$)}
\sms

\section{Logical expressions}
Relational operators yield values {\tt .TRUE.} or {\tt .FALSE.}.
The operator takes two forms, a word with period before and after it
(Fortran 77), or C-like notation (Fortran 90).
\sms
\halign{\t\tt#\hfil\t&\t\tt#\hfil\t&\t#\hfil\cr
.LT. & < & less than \cr
.LE. & <= & less or equal \cr
.EQ. & == & equal to \cr
.NE. & /= & not equal \cr
.GT. & > & greater than \cr
.GE. & >= & greater or equal \cr
}
\sms
The logical operators in order of decreasing precedence are
\sms
\vbox{
\halign{\t\tt#\hfil\t&\t#\hfil\cr
.NOT. &  negation \cr
.AND. &  and/intersection \cr
.OR. &  or/union \cr
.EQV. {\rm and} .NEQV. &  equivalence or non-equivalence \cr
}}
\sms
Here are some examples
\ttbg
I = .FALSE.
J = .TRUE.
L = X < 0
K = I .OR. J .AND. .NOT. L
K = N > 0 .AND. N < 10
\tted

\section{Character expressions}
Character strings are quoted by apostrophes.  Fortran~90 also uses
double quotes. E.g.,
\sms
\ll{'a string'\h 'a single quote '' inside'\h "a single quote ' inside"}
\sms
To produce a single quote when the string is quoted with single quotes, 
one needs two single quotes.  Of course, one can avoid this inconvenience 
by quoting the string in double quotes.  The only operator for
character string is the concatenation operator {\tt //}, e.g., 
{\tt X // '.f90'}, {\tt 'AB' // 'CD'} produces {\tt 'ABCD'}.  Other
examples involving character expressions:
\sms
\l{\v`CHARACTER(LEN=4) CHAR1, CHAR2`\m{declare four letter variables}}
\l{\v`CHARACTER(LEN=8) RESULT`\m{{\ninett RESULT} can hold up to 8 
                                  characters}}
\l{\v`CHAR1 = 'ANY '`\m{{\ninett CHAR1} gets {\ninett 'ANY '}}}
\l{\v`CHAR2 = 'books'`\m{{\ninett CHAR2} gets {\ninett 'book'}}}
\l{\v`RESULT = CHAR1 // CHAR2`\m{{\ninett RESULT} has {\ninett 'ANY book'}}}
\sms
You can refer to each individual character or substring with the
$i${\tt :}$j$ notation. Continue the previous examples,
\sms
\l{\v`CHAR1(1:2) = 'IV'`\m{{\ninett CHAR1} becomes {\ninett 'IVY '}}}
\l{\v`RESULT(6:7) = RESULT(1:2)`\m{{\ninett RESULT} now has {\ninett 'ANY bANk'}}}
\l{\v` `\m{{\ninett RESULT(1:1)} has {\ninett 'A'}}}
\l{\v` `\m{{\ninett RESULT(5:8)} has {\ninett 'bANk'}}}
\sms

\section{Operator precedence}
Here is a complete list of operators and their precedence in 
decreasing order
\sms
\halign{\t\tt\hfil#\hfil\t&\t\hfil#\hfil\t\cr
{\rm Operator}& Precedence \cr
{\rm user-defined unary operator}& highest \cr
**& $\Downarrow$ \cr
*{\rm,\quad}/& \cr
{\rm (unary)} +{\rm,\quad}-& \cr
{\rm (binary)} +{\rm,\quad}-& \cr
.EQ.{\rm,} .NE.{\rm,} .LT.{\rm,} .LE.{\rm,} .GT.{\rm,} .GE.& \cr
=={\rm,} /={\rm,} <{\rm,} <={\rm,} >{\rm,} >=& \cr
.NOT.& \cr
.AND.& \cr
.OR.& \cr
.EQV.{\rm,} .NEQV.& $\Downarrow$ \cr
{\rm user-defined binary operator}&lowest \cr
}
\sms
Associativity is from left to right, with the exception of exponentiation
operator {\tt **}, which is from right to left.

\section{structured data types}
Fortran,  especially Fortran~90, has a rich collections of operations
and constructors on arrays.  For example, an array can be assigned to
another array in one statement
\sms
\ll{B = A}
\sms
We'll discuss them later.  Fortran~77 does not have structures
and pointers.  They are Fortran~90 features.

Here is another simple program
\ttbg
PROGRAM EXAMPLE
! Calculate the area of a circle
REAL, PARAMETER :: PI = 3.1415926
REAL :: RADIUS, AREA
CHARACTER(LEN=16) :: CHAR

CHAR = 'Area of a circle'
READ *, RADIUS
IF (RADIUS < 0) STOP 'Radius < 0'
AREA = PI * RADIUS ** 2
PRINT *, CHAR // ' is', AREA
END PROGRAM EXAMPLE
\tted


\vfill\eject
\lecture{3. Control Statements}

\section{IF construct}
A conditional if in Fortran takes the form
\sms
\l{\t\tt IF ($logical$-$expr$) THEN\m{if logical expression yields value
{\ninett .TRUE.},}}
\l{\t\tt\ \ \ $block$\m{the block (any number of lines) is executed}}
\l{\t\tt END IF\m{or equivalently {\ninett ENDIF}}}
\sms
If the $block$ contains one statement, one can simply write
\sms
\l{\t\tt IF ($logical$-$expr$) $statement$\m{single line}}
\sms
One can insert any number of {\tt ELSE IF} and one {\tt ELSE} before
the end as
\sms
\ll{IF ($expr_1$) THEN}
\ll{\ \ \ $block_1$}
\ll{ELSE IF ($expr_2$) THEN}
\ll{\ \ \ $block_2$}
\ll{.....}
\ll{ELSE}
\ll{\ \ \ $block_n$}
\ll{\tt END IF}
\sms
The first $block_i$ with $expr_i$ yields {\tt .TRUE.} is executed.  If none
of the $expr$ is true, {\tt ELSE} block is executed. E.g.,
\ttbg
SWAP: IF (X < Y) THEN
         TEMP = X
         X = Y
         Y = TEMP
      END IF SWAP
\tted
Here the word {\tt SWAP} is optional, and is to help the reader to 
see the if block more clearly.  
\ttbg
IF (I < 0) THEN
   IF (J < 0) THEN
      X = 0.
      Y = 0.
   ELSE
      Z = 0.
   END IF
ELSE IF (K < 0) THEN
   Z = 1.
ELSE
   X = 1.
   Y = 1.
END IF
\tted

\section{CASE construct}
The {\tt CASE} construct is not available in Fortran 77, it is
similar to Pascal's:
\sms
\l{\v`SELECT CASE (NUMBER)`\m{{\ninett NUMBER} is an integer type}}
\l{\v`   CASE (:-1)`\m{all values of {\ninett NUMBER <= -1}}}
\l{\v`      NSIGN = -1`\m{}}
\l{\v`   CASE (0)`\m{{\ninett NUMBER} exactly zero}}
\l{\v`      NSIGN = 0`\m{}}
\l{\v`   CASE (1:4)`\m{values between 1 and 4 inclusive}}
\l{\v`      NSIGN = 2`\m{}}
\l{\v`   CASE (5:)`\m{all values 5 or above}}
\l{\v`      NSIGN = 1`\m{}}
\l{\v`END SELECT`\m{}}
\sms
The general form of case selector can be a list of non-overlapping
values and range, separated by comma, e.g., 
\sms
\ll{CASE (1, 7, 10:13, 23)}
\sms
In the above example, 
the case is selected if the value is 1, 7, 10, 11, 12, 13, and 23. 
{\tt CASE(:0)} selects all values smaller than or equal to zero.
The variables can be character and logical values besides
integers.  E.g.,
\sms
\l{\v`SELECT CASE (C)`\m{{\ninett C} is character type}}
\l{\v`   CASE ('C')`\m{}}
\l{\v`      CHAR_TYPE = .TRUE.`\m{}}
\l{\v`   CASE ('I':'N')`\m{}}
\l{\v`      INT_TYPE = .TRUE.`\m{}}
\l{\v`   CASE DEFAULT`\m{any other case not listed above}}
\l{\v`      REAL_TYPE = .TRUE.`\m{}}
\l{\v`END SELECT`\m{}}
\sms

\section{DO construct}

DO-loop is probably the most important and most
frequently used construct.
\sms
\ll{DO $variable$ = $start$, $end$, $step$}
\ll{\ \ \ $block$}
\ll{END DO}
\sms
where $start$ and $end$ are the starting and ending value of the loop
$variable$.  The $variable$ is incremented in $step$, if omitted,
default to 1.  The step can be positive or negative.  The loop will 
iterate exactly 
$\max\bigl( (end - start + step)/step, 0\bigr)$ times.
E.g.,
\sms
\l{\v`SUM = 0.0`\m{initialize the {\ninett SUM}}}
\l{\v`DO I = 1, 10`\m{index {\ninett I} varies from 1 to 10 in step 1}}
\l{\v`   SUM = SUM + A(I)`\m{{\ninett A(I)} is array element}}
\l{\v`END DO`\m{}}
\sms
Fortran 77 style is 
\sms
\l{\v`      DO 200 I = 1, 10`\m{200 is a label for the line ending the loop}}
\l{\v`200   SUM = SUM + A(I)`\m{loop end here, {\ninett END DO} not allowed}}
\sms
Fortran 90 has a more general form of {\tt DO} as
\sms
\l{\v`DO`\m{infinite loop results if {\ninett EXIT} does not appear in block}}
\ll{\ \ \ $block$}
\ll{END DO}
\sms
A statement {\tt IF($something$) EXIT} must be in the $block$ in
order to terminate the loop.   The statement 
\sms
\ll{IF ($something$) CYCLE}
\sms
transfers the control immediately to the {\tt END DO} statement and 
continue for the next iteration if needed.

Fortran~90 also has a 
\sms
\ll{DO WHILE (...)}
\ll{\ \ \ ...}
\ll{END DO}
\sms
You can get the same effect with 
\sms
\ll{DO}
\ll{\ \ \ IF (...) EXIT}
\ll{\ \ \ ...}
\ll{END DO}
\sms
Pascal's repeat-until can be emulated by putting the exit condition 
at the end.

\section{GO TO statement}
{\tt GO TO} directs the program to a labelled statement, e.g.,
\ttbg
      X = Y + 3
      IF (X .GT. 1) GO TO 10
      X = X + 2
10    Z = X + Y
\tted
A label is a number occupying column one to five.   It is illegal
to jump into a {\tt DO} or other blocks from outside with {\tt GO TO}.

\lecture{4. Basic Program Units and Procedures}

\section{Main program}

A main program in Fortran takes the form
\sms
\l{\v`PROGRAM ILLUSTRATE`\m{the {\ninett PROGRAM} statement is optional}}
\l{\v`   REAL :: X, Y`\m{variable declaration}}
\l{\v`   PRINT *, 'Enter x and y'`\m{}}
\l{\v`   READ *, X, Y`\m{read in data}}
\l{\v`   PRINT *, 'Squares of the numbers are', SQ(X), SQ(Y)`\m{function call}}
\l{\v`   CALL SWAP(X, Y)`\m{call a subroutine}}
\l{\v`   PRINT *, 'X = ', X, ' Y =', Y`\m{}}
\l{\v`END PROGRAM ILLUSTRATE`\m{the words {\ninett PROGRAM ILLUSTRATE} 
				are optional}}
\sms

\section{External functions and subroutines}

The above program calls one function {\tt SQ()} and a subroutine 
{\tt SWAP}.  An external function or subroutine can appear in another
file,  or in the same file before or after the main program. 
We'll adopt the convention of putting the subprogram after the main
one.   The function definition may look like this
\sms
\l{\v`REAL FUNCTION SQ(X)`\m{define a function}}
\l{\v`   REAL :: X`\m{we'll indent three spaces every block structure}}
\l{\v`   SQ = X * X`\m{{\ninett SQ} must appear on the lefthand side
at least one}}
\l{\v`END FUNCTION SQ`\m{the words {\ninett FUNCTION SQ} are optional}}
\sms
A subroutine does not return any value, here is the swap function
\sms
\l{\v`SUBROUTINE SWAP(A, B)`\m{define a subroutine}}
\l{\v`   REAL :: A, B, TMP`\m{declaration of arguments and local variable}}
\l{\v`   TMP = A`\m{}}
\l{\v`   A = B`\m{}}
\l{\v`   B = TMP`\m{}}
\l{\v`END SUBROUTINE SWAP`\m{Fortran 77 does not allow the optional
words {\ninett SUBROUTINE SWAP}}}
\sms
In Fortran arguments are always passed by reference, just the opposite of 
C, and like Pascal's {\tt var} parameters. 

\section{Internal functions and subroutine}
This feature is available only in Fortran 90.  A (main) program,
an external function, an external subroutine, or a module can contain
inside it internal functions and/or subroutines.   
Unlike Pascal, this nesting is one-level only.  Internal 
subprograms cannot contain further subprograms.  Variables declared
in the  main program (or external function/\-subroutine) are visible
to the internal subprogram.  Here is an example containing internal
subprograms.  
\ttbg
PROGRAM SORT
   IMPLICIT NONE
   REAL :: N1, N2, N3
\tted
\ttbg
   CALL READ
   CALL SORT
   CALL PRINT
\tted
\sms
\l{\v`CONTAINS`\m{the internal subroutines follows}}
\sms
\ttbg   
   SUBROUTINE READ
      READ *, N1, N2, N3
   END 

   SUBROUTINE SORT
      IF (N1 > N2) CALL SWAP(N1, N2)
      IF (N1 > N3) CALL SWAP(N1, N3)
      IF (N2 > N3) CALL SWAP(N2, N3)
   END

   SUBROUTINE PRINT
      PRINT *, N1, N2, N3
   END 

END PROGRAM SORT
\tted     

\section{Keyword and optional arguments}
This function calculates the sum $\sum_{i = n}^m x^i$,
\ttbg
FUNCTION SERIES(X, LOWER, UPPER)
   INTEGER :: I, LOWER, UPPER
   REAL :: SERIES, X
   SERIES = 0.0
   DO I = LOWER, UPPER
      SERIES = SERIES + X**I
   END DO
END FUNCTION SERIES
\tted
The usual way of evoking the function is, say, {\tt SERIES(0.5, 1, 10)}, 
where {\tt X} gets the value 0.5, {\tt LOWER} and {\tt UPPER} have
1 and 10, respectively.  But it can also be called as
\sms
\l{\t\tt SERIES(0.5,LOWER=1,UPPER=10) \hfill {\rm or}
\hfill SERIES(0.5,1,UPPER=10) \hfill {\rm or}  \hfill
SERIES(UPPER=10,LOWER=1,X=0.5)\t}
\sms
We can even omit some of the arguments if we declare them optional,
rewriting the  above function as 
\ttbg
FUNCTION SERIES(X, LOWER, UPPER)
   INTEGER :: I, ISTART, UPPER
   REAL :: SERIES, X
   INTEGER, OPTIONAL :: LOWER
   
   IF ( PRESENT(LOWER) ) THEN
      ISTART = LOWER
   ELSE
      ISTART = 1
   END IF

   SERIES = 0.0
   DO I = ISTART, UPPER
      SERIES = SERIES + X**I
   END DO
END FUNCTION SERIES
\tted
then we can call the function as
\sms
\ll{SERIES(0.5, LOWER = 1, UPPER = 10) \h {\rm or} \h
SERIES(0.5, UPPER = 10)}
\sms
The intrinsic function {\tt PRESENT()} detects the presence of the 
optional argument.  If it is absent from the argument list, 1 is
used as default. 

\section{Interface}
Interface is somewhat like function prototypes in C.  It  aids compiler
to generate correct codes, as well as does some other things.  
In the above optional argument example, if the 
function is used as an external function, then an interface block (in the
declaration part) is required in the calling program:
\ttbg
INTERFACE
   FUNCTION SERIES(X, LOWER, UPPER)
      INTEGER :: UPPER
      INTEGER, OPTIONAL :: LOWER
      REAL :: SERIES, X
   END FUNCTION SERIES
END INTERFACE
\tted
In an interface block, the name of dummy arguments can be renamed.  Those
renamed argument then can be used as keywords.

Interface block can also be used to define a generic procedure 
(function or subroutine).  A generic function can 
accept several different data types as its arguments.  For example, 
{\tt SIN(X)} works for both {\tt X} being a real number as well as
complex number. {\tt SIN} is a generic name for the sine function.
This is achieved by written several functions for each type, and let the
compiler choose the appropriate  one.  Here is an example for a generic
{\tt SWAP} function.  First we define two specific functions, one for
integer swap, and the other for real number.
\sms
\l{\v`SUBROUTINE SWAP_INTEGER(A, B)`\m{integer version}}
\l{\v`   INTEGER :: A, B, TMP`\m{}}
\l{\v`   TMP = A`\m{}}
\l{\v`   A = B`\m{}}
\l{\v`   B = TMP`\m{}}
\l{\v`END SUBROUTINE SWAP_INTEGER`\m{}}
\sms
\sms
\l{\v`SUBROUTINE SWAP_REAL(A, B)`\m{real value version}}
\l{\v`   REAL :: A, B, TMP`\m{}}
\l{\v`   TMP = A`\m{}}
\l{\v`   A = B`\m{}}
\l{\v`   B = TMP`\m{}}
\l{\v`END SUBROUTINE SWAP_REAL`\m{}}
\sms
An interface block specifies a generic name {\tt SWAP}
\ttbg
INTERFACE SWAP
   SUBROUTINE SWAP_INTEGER(A, B)
      INTEGER :: A, B
   END SUBROUTINE SWAP_INTEGER

   SUBROUTINE SWAP_REAL(A, B)
      REAL :: A, B
   END SUBROUTINE SWAP_REAL
END INTERFACE
\tted
After this, one can use {\tt SWAP(A, B)} for the arguments both integers
or reals.

\section{Procedures as arguments}
In Fortran, we can also pass function or subroutine name as an argument,
as illustrated in the following example:
\ttbg
PROGRAM MAIN
   REAL A, B, F
   INTERFACE
      REAL FUNCTION FUN(X)
         REAL :: X
      END FUNCTION
   END INTERFACE
   F = MINIMUM(1.0, 2.0, FUN)
   ...
END PROGRAM

REAL FUNCTION MINIMUM(A, B, FUNC)
   REAL :: A, B, F, X
   INTERFACE
      REAL FUNCTION FUNC(X)
         REAL :: X
      END FUNCTION FUNC
   END INTERFACE
   ...
   F = FUNC(X)
   ...
END FUNCTION MINIMUM

REAL FUNCTION FUN(X)
   ...
\tted

\lecture{5. Input, Output, and Files}

\section{Formatted input and output}

The input/output format specification is rather sophisticated in Fortran. 
We will just introduce the most common ones.   Here is some examples of
formatted outputs (assuming that {\tt I} is an integer, {\tt A} is an
array of size 10, and {\tt WORD} is a string of character of length 20):
\sms
\l{\v`PRINT *, I, A(1)`\m{use the default format}}
\l{\v`PRINT '(I10)', I`\m{integer in a field width of 10}}
\l{\v`PRINT '("I =", I10)', I`\m{character string can be in the 
format specification}}
\l{\v`PRINT "(10F10.3)", A`\m{10 reals of width 10, precision 3, 
{\ninett "} and {\ninett '} are equivalent}}
\l{\v`PRINT '(5F10.3, E10.2)', (A(I), I = 1, 9, 2), A(2)`\m{implicit loop}}
\l{\v`PRINT '(A10)', WORD(5:14)`\m{five characters from fifth to fourteenth}}
\l{\v`PRINT 100, A(1), A(2)`\m{use {\ninett FORMAT} statement}}
\l{\v`100 FORMAT(1X, 2F10.4)`\m{format for the above output}}
\sms
Input with {\tt READ} is similar.  In general the basic I/O takes the form:
\sms
\ll{PRINT $fmt$, $variable$ $list$}
\ll{READ $fmt$, $variable$ $list$}
\sms
The format $fmt$ is {\tt *} (default format), or a number
(the format is specified by a corresponding {\tt FORMAT} statement), or
a string $fmt$ of the form
\sms
\ll{'($r$D$w$,...)'}
\sms
where $r$ is repeat factor which indicates that edit descriptor that
follows will be repeated $r$ times.  {\tt D} specifies the type of
format, which can be, among many others, one of the following
\sms
\halign{\t\tt#\hfil\h\h&#\hfil\cr
I$w$ & integer with field width $w$\cr
F$w$.$d$ & real number of width $w$, $d$ digits after decimal point\cr
A {\rm\quad or\quad} A$w$ & character string (width $w$)\cr
L$w$ & logical data, printed as {\tt T} or {\tt F}\cr
X & space\cr
T$n$ & tab to position $n$\cr
/ & new line\cr
B$w$\ \ O$w$\ \ Z$w$ & integer in binary, octal, hexadecimal notation\cr
E$w$.$d$ & exponential notation for real number\cr
G$w$.$d$ & use {\tt F} or {\tt E} format, whichever is shorter\cr
}
\sms
Here is an example that uses last three format descriptors:
\ttbg
I = 10; J = 12; K = 5; R = 3.1415; W = 'word'
PRINT '(T3, I2, 2X, I2, /, I3)', I, J, K
PRINT '(F5.2, E9.2)', R, R
PRINT '(B5)', K
PRINT '(A)', W
PRINT '(I1, F3.1)', J, 1001.12
\tted
will print out
\sms
\ll{\char`\ \char`\ \char`\ 10\char`\ \char`\  12}
\ll{\char`\ \char`\ 5}
\ll{\char`\ 3.14\char`\ 0.31E+01}
\ll{\char`\ \char`\ 101}
\ll{word}
\ll{****}
\sms
Format statement can be shared by different I/O statements, e.g.,
\ttbg
PRINT 100, I, J
PRINT 100, J, K, R
100 FORMAT(1X, I10, 2X, I10, 2X, F15.5)
\tted

\section{Operation on files}
To access file, we must open the file with, e.g.,
\ttbg
OPEN(2, FILE = 'F.DAT', STATUS='NEW')
\tted 
The first number is called {\tt UNIT} number (between 1 to 99) 
assigned to the file.   Then the file name to be opened.  {\tt STATUS} 
should be {\tt 'NEW'} (new file), {\tt 'OLD'} (existing file), or 
{\tt 'UNKNOWN'} and some other options.   

File should be closed, e.g., 
\sms
\l{\v`CLOSE(2)`\m{equivalent to {\ninett CLOSE(UNIT = 2)}}}
\sms
Read/write operations to/from the file are with the statements
\sms
\ll{READ(UNIT = $u$, FMT = $fmt$) $list$}
\ll{WRITE(UNIT = $u$, FMT = $fmt$) $list$}
\sms
where $u$ is the unit number.  The keyword {\tt UNIT} is optional if
the unit number is the first argument.  And  $fmt$ is the format
specification.   Also, the keyword {\tt FMT} is optional if 
$fmt$ takes the second argument. 

There are quite a number of other optional keyword specifications
for {\tt OPEN}, {\tt WRITE}, and {\tt READ}.   The most
useful one probably is {\tt ERR = $label$}.  If an error occurs
during operation, control transfers to the statement with the integer
$label$.  The other keyword is  {\tt IOSTAT = $integer$-$variable$}.  
The variable gets, (1) a positive value if an input error occurs, (2)
a negative value if end-of-file is encountered, (3) and zero
otherwise. 

The {\tt PRINT} statement is equivalent to {\tt WRITE(*,...)},
namely, write to the standard output.  When a unit number is an
asterisk {\tt *}, the standard I/O is used.   On most systems, the
unit number 5 is reserved for standard input (keyboard), and 
unit number 6 is for standard output (screen).  These numbers should
be avoided when open a file.  Here are some more examples:
\ttbg
    READ(5,*) A, B, C
    WRITE(*, "(E20.5)") X
    WRITE(2, 102) X, Y
    ...
102 FORMAT('X=', F10.4, ' Y=', F10.4)
\tted

\section{Internal files}
The external files are associated with the external disc storage. 
Internal files are just text in main CPU memory.  For internal file
I/O, the unit number is replaced by the name of character string --- the
string is the file.  E.g.
\sms
\ll{CHARACTER (LEN = 120) :: BUFFER}
\l{\v`WRITE (BUFFER, 'I10') 101`\m{write number to the string}}
\l{\v`READ (BUFFER, *) I`\m{read from the string}}
\sms
Using internal file is an effect way of converting numbers to 
character strings.

\section{Unformatted I/O}
Formatted I/O converts numbers into strings of ASCII characters.  Unformatted
I/O is a direct representation without conversion.  If a large amount of
data needs to be written, this unformatted I/O is much more efficient.
however, since the data are in binary, they can not be read directly
by human.  Here is an example of using unformatted output.
\ttbg
OPEN(9, FILE ='TMP', FORM='UNFORMATTED')
WRITE(9) X, Y, Z
WRITE(9, ERR=199) A, B, C
CLOSE(9)
\tted
Since the output is unformatted, you need not supply a format 
specification.

\section{Sequential vs. direct access}
Fortran considers each line as a record.  Sequential access read/write
records in sequence.  Direct access read/write records in arbitrary order.
E.g.,
\sms
\l{\v`OPEN(7, FILE='RECORD', ACCESS='DIRECT')`\m{open for direct access}}
\l{\v`READ(7,*,REC=10) A`\m{read line 10}}
\l{\v`READ(7,*,REC=3) B`\m{read line 3}}
\l{\v`READ(7,*,REC=17) C`\m{read line 17}}
\l{\v`CLOSE(7)`\m{close the file}}
\sms

\lecture{6. Arrays}

\section{Array declarations}
There are a number of ways to declare arrays 
\sms
\l{\v`INTEGER A(10)`\m{with elements {\ninett A(1)} to {\ninett A(10)}, 
					      Fortran 77 style}}
\l{\v`CHARACTER CH(0:3)`\m{declare {\ninett CH(0)}, {\ninett CH(1)}, 
				{\ninett CH(2)}, {\ninett CH(3)}}}
\l{\v`INTEGER, DIMENSION(-9:9) :: B`\m{with elements {\ninett B(-9)} to 
			          {\ninett B(9)}}}
\l{\v`REAL, DIMENSION(10, 10) :: C`\m{a 10 by 10 matrix}}
\l{\v`INTEGER, DIMENSION(:), ALLOCATABLE ::D`\m{memory allocated dynamically}}
\sms

\section{Allocatable arrays} 

Note that in the last declaration, {\tt D} is specified as
one-\-dimensional, but the exact array bounds are not given, and it has
the attribute {\tt ALLOCATABLE}.   The actual bounds and memory for it can
be specified at run time.
The dynamic storage allocation  is not possible in Fortran 77.
With
\sms
\ll{ALLOCATE(D(1:10))}
\sms
memory is allocated when running the program.
Now {\tt D} is an array of size 10.  After use, the memory can be freed with
\sms
\ll{DEALLOCATE(D)}
\sms

\section{Array constructors}
The features discussed in the rest of this lecture are available only in 
Fortran 90. If {\tt X}, {\tt B}, {\tt C}, and {\tt D} are one-dimensional 
arrays of real, size 4, then we can assign value to each of the elements 
in a single statement.
\sms
\l{\v`X = (/ 1.2, 2.0, 0.0, -3.5 /)`\m{four individual values}}
\l{\v`X = (/ A(I, 1:2), A(I+1, 2:3) /)`\m{array section is used}}
\l{\v`X = (/ ( SQRT(REAL(I)), I = 1, 4 ) /)`\m{implied do loop is used}}
\l{\v`X = 1.0`\m{every element gets 1.0}}
\l{\v`X = SQRT(D)`\m{applies intrinsic function to array}}
\l{\v`X = A(I, 1:4) * D + B/C`\m{corresponding element is computed}}
\sms

The last three examples are not array constructors, but array
expressions and assignment. Operators and functions may be applied to
arrays having the same number of elements (same shape).  The operation is 
carried out elementwise.  The last example above is equivalent to
\ttbg
DO J = 1, 4
   X(J) = A(I,J) * D(J) + B(J)/C(J)
END DO
\tted
\sms
\section{Array sections}
Array names like {\tt X}, {\tt A}, {\tt B}, $\ldots$, refer to the
whole arrays.  In C they refer only to the starting addresses of the
arrays. In the second example of array constructor we use expressions like
{\tt A(I, 2:3)},  This is a one-dimensional array with elements
{\tt A(I, 2)} and {\tt A(I, 3)}.  It is called an array section. 
In general if {\tt X} is an array of certain size, then 
\sms
\ll{X(n:m)}
\sms
is the array with elements 
\sms
\ll{X(n), X(n+1), ... X(m)}
\sms
Beside the lower and upper bounds, one can also specify a {\sl
stride} as, e.g.,  
\sms
\ll{X(3:7:2)}
\sms
It represents an array with three elements
\sms
\ll{X(3), X(5), X(7)}
and 
\ll{X(7:1:-1)}
\sms
is an array with elements
\sms
\ll{X(7), X(6), X(5), ..., X(1)}
\sms


The indices of an array can be another integer array (vector), for
example,  if {\tt INDEX = (/ 1, 4, 2, 2 /)}, then {\tt X(INDEX)} 
represents an array section with elements
\sms
\ll{X(1), X(4), X(2), X(2)}
\sms
Other examples of array section are 
\sms
\l{\v`A(I, :)`\m{the whole row {\ninett I} of array {\ninett A}}}
\l{\v`A(:, J)`\m{the whole column {\ninett J} of array {\ninett A}}}
\l{\v`A(I, 10:1:-1)`\m{the row {\ninett I} in reverse order}}
\l{\v`A(I, :5)`\m{the row {\ninett I} from lower bound to 5}}
\l{\v`A(2:, J)`\m{the column {\ninett J} from 2 to upper bound}}
\sms
Array sections can be used just like whole arrays in expressions or 
assignments as long as they have the same size and shape, e.g.,
\sms
\ll{A(2:4, 5:8) = A(3:5, 1:4)}
\ll{X(1:4) = A(1, 2:5)}
\sms

\section{Where construct}
{\tt WHERE} construct is the array equivalent of the {\tt IF}
construct for scalar quantities.  Here is an example
\ttbg
WHERE (C /= 0)
   A = B / C
ELSEWHERE
   A = 0
   C = 1
END WHERE
\tted
All of the array {\tt A}, {\tt B}, and {\tt C} must have the same
size and shape.  If not, array sections must be used.
The operation is on an element-\-by-\-element basis.
If the arrays have shape or dimension (10, 10), 
the above statement is equivalent to 
\ttbg
DO I = 1, 10
   DO J = 1, 10
      IF (C(I,J) /= 0) THEN
         A(I,J) = B(I,J)/C(I,J)
      ELSE
         A(I,J) = 0
         C(I,J) = 1
      END IF
   END DO
END DO
\tted
Of course, the {\tt WHERE} construct is much shorter and elegant.

\section{Array intrinsic functions}
Many intrinsic functions can be applied to array.  The result is an
array applying the function elementwise. E.g.,
\ttbg
PI = 3.14159265358979
A = (/0, PI/2, PI, 2*PI /)
\tted
then {\tt COS(A)} is an array with elements $\cos 0$, $\cos \pi/2$,
$\cos\pi$, and $\cos 2\pi$, or {\tt (/ 1, 0, -1, 1 /)}.

There is a large collection of intrinsic functions for arrays, 
here are most of them:
\sms
\ll{SUM(A)\hfill $\sum_i A(i)$\t}
\ll{PRODUCT(A)\hfill $\prod_i A(i)$\t}
\ll{MAXVAL(A)\hfill $\max_i A(i)$\t}
\ll{MINVAL(A)\hfill $\min_i A(i)$\t}
\ll{DOT\_PRODUCT(A, B)\hfill $\sum_i A(i) B(i)$\t}
\ll{MATMUL(A, B)\hfill $\sum_k A(i,k) B(k,j)$\t}
\ll{TRANSPOSE(A)\hfill\rm matrix transpose\t}
\ll{SHAPE(A)\hfill\rm A 1-dimensional array, (2,5,5) means a 
$2 \times 5 \times 5$ array\t}
\ll{SIZE(A)\hfill\rm total number of array elements\t}
\ll{MAXLOC(A)\hfill\rm location of maximum\t}
\ll{MINLOC(A)\hfill\rm location of minimum\t}
\ll{ALL(M)\hfill\rm {\tt M} is logical variable, return {\tt .TRUE.}
if all elements is {\tt .TRUE.}\t}
\ll{ANY(M)\hfill\rm return {\tt .TRUE.} if any element is {\tt .TRUE.}\t}
\sms
There are a number of other array manipulation functions
{\tt RESHAPE} changes the shape of an array, e.g.,
\sms
\ll{RESHAPE(SOURCE = (/1, 2, 3, 4, 5, 6/), SHAPE = (/2, 3/))}
\sms
has the value (a $2\times 3$ array)
\ttbg
1  3  5
2  4  6
\tted
{\tt SPREAD} enlarges and duplicates array,
\ttbg
SPREAD(SOURCE = (/1, 2, 3/), DIM = 1, NCOPIES=3)
\tted
has the value (a $3\times 3$ matrix)
\ttbg
1  2  3
1  2  3
1  2  3
\tted
{\tt CSHIFT} circularly shifts the array, 
\ttbg
CSHIFT(ARRAY = (/1, 2, 3/), SHIFT = 1)
\tted
has the value (a one dimensional array)
\ttbg
2  3  1
\tted
{\tt PACK} picks out elements and put them in one dimension. E.g., if
{\tt A} is left integer array and {\tt M} is the right logical array,
\ttbg
1  2  3              T  F  T
4  5  6              T  T  F
7  8  9              F  F  T
\tted
then 
\sms
\ll{PACK(A, M)}
\sms
has the value (a one dimensional array)
\sms
\ll{(/1, 4, 5, 3, 9/)}
\sms

\section{Zero-sized arrays}
{\tt A(3:2)} is a zero-sized array which contains no elements.  Zero-size
array causes program ``do nothing.''  E.g.,
\ttbg
DO I = 1, N
   X(I) = B(I)/A(I,I)
   B(I+1:N) = B(I+1:N) - A(I+1:N,I) * X(I)
END DO
\tted
The above code solves the lower-triangular set of linear equations.
$$\eqalign{
a_{11} x_1 &= b_1,\cr
a_{21} x_1 + a_{22} x_2 & = b_2,\cr
\cdots  & \cr
a_{n1} x_1 + a_{n2} x_2  + \cdots a_{nn} x_n & = b_n.\cr
}$$
We encounter a zero-sized array when the index {\tt I} takes its final
value {\tt N}.

\section{Arrays of assumed shape}
In a subroutine or function where array is passed as an argument, 
the actual bounds need not declared.  It is assumed from the calling
program.  When this feature is used, an explicit interface is
required.  E.g.
\ttbg
SUBROUTINE SUB(A)
   INTEGER, DIMENSION(:,:,:) :: A
   ...
\tted

\section{Automatic arrays}
In the following example, the array {\tt TMP} is called automatic array.
Automatic array must be a local variable, and can not be a dummy
argument.
\ttbg
SUBROUTINE SWAP(A,B)
   REAL, DIMENSION(:) :: A, B
   REAL, DIMENSION(SIZE(A)) :: TMP
   TMP = A
   A = B
   B = TMP
END SUBROUTINE
\tted

\section{Array-valued functions}
Fortran 90 allows a function to return an array.   The following function
multiplies a vector {\tt B} by a matrix {\tt A}, and returns the resulting
vector.
\ttbg
FUNCTION MULT(A, B)
   REAL, DIMENSION(:,:) :: A
   REAL, DIMENSION(SIZE(A,2)) :: B
   REAL, DIMENSION(SIZE(A,1)) :: MULT
   INTEGER :: J, N
   N = SIZE(A,1)
   MULT = 0.
   DO J = 1, SIZE(A,2)
      MULT = MULT + A(1:N, J) * B(J)
   END DO
END FUNCTION
\tted
The function can be used as, e.g.,
\ttbg
C(1:10) = MULT(A(1:10, 1:5), B(1:5))
\tted
Interface for such function must be explicit, i.e., if the function is
an external function, an interface specification is required.  Interface
is not needed for internal functions (functions defined after the
keyword {\tt CONTAINS}).


\lecture{7. Character Data}

\section{Declarations}
There is no distinction between string and character, a single character
is a string of length one.  The character variable declaration takes the
form
\sms
\l{\v`CHARACTER (LEN = 8) :: STRING_8`\m{defined a string of length 8}}
\l{\v`CHARACTER*5 NAME`\m{character string of length 5, Fortran 77 style}}
\l{\v`CHARACTER :: LETTER`\m{a single character}}
\l{\v`CHARACTER (LEN = 80), DIMENSION(60) :: PAGE`\m{array of `strings'}}
\l{\v`CHARACTER*15 ACH(100)`\m{Fortran 77 style}}
\sms
In a subroutine or function, the length of the character string can be
declared as {\tt LEN = *}.  The actual length is assumed from the calling
program. 

\section{Assignment, operations, and substrings}
The character string values are quoted in apostrophes or double quotes,
e.g.,
\sms
\l{\v`STRING_8 = 'a string'`\m{Fortran 90 also accepts double quotes}}
\l{\v`NAME = 'It''s '`\m{two apostrophes produce one, one can say 
{\ninett "It's\char`\ "}}}
\l{\v`LETTER = 'letter'`\m{{\ninett 'l'} is actually assigned}}
\l{\v`ACH(2) = "array"`\m{{\ninett "array\char`\ \char`\ \char`\ \char`\ \char`\ \char`\ \char`\ \char`\ \char`\ \char`\ "} is assigned}}
\l{\v`PAGE(1) = "a long line&`\m{continuation of a string}}
\l{\v`          & continued&`\m{}}
\l{\v`          & with ampersand & is okay."`\m{}}
\sms
If the string on the righthand side is shorter than the declared variable
length, it is appended with blanks.  If it is longer, the tail part is
truncated.  Two character strings can be concatenated 
\sms
\l{\v`NAME//STRING_8`\m{is a string of length 13, having value 
{\ninett It's a string}}}
\sms
Beside the concatenation operator {\tt //}, 
the comparison operators like {\tt ==} or {\tt <} can be applied to strings.
E.g.,
\sms
\l{\v`NAME == "It's "`\m{has value {\ninett .TRUE.}}}
\l{\v`NAME /= "This"`\m{has value {\ninett .TRUE.}}}
\l{\v`NAME > STRING_8`\m{has value {\ninett .FALSE.}}}
\l{\v`"12" > "2"`\m{is {\ninett .FALSE.}}}
\l{\v`"apple" < "pear"`\m{is {\ninett .TRUE.}}}
\sms
A substring can be referred with range indices, similar to array section.
For example,
\sms
\l{\v`STRING_8(3:8)`\m{has value {\ninett string}}}
\l{\v`NAME(4:4)`\m{has value {\ninett s}}}
\l{\v`NAME(3:3) = 'o'`\m{changes {\ninett "It's\char`\ "} to {\ninett "Itos\char`\ "}}}
\l{\v`PAGE(1)(3:6)`\m{has value {\ninett "long"}}}
\sms
Substring can be used just like string, they can appear on the right or
left side of an assignment. 


\section{Character intrinsic functions}
There are a number of intrinsic functions on character strings.
\sms
\halign{\t\tt#\hfil\h&#\hfil\cr
LEN($string$) & length of the $string$, including trailing blanks\cr
LEN\_TRIM($string$) & length of the $string$, ignoring trailing blanks\cr
TRIM($string$) & string with trailing blanks removed\cr
INDEX($string$, $sub$) & the position of $sub$ string in $string$\cr
CHAR($code$) & the character with integer $code$\cr
ICHAR($c$) & the code of the character $c$\cr
ADJUSTL($string$) & adjust left of the string\cr
ADJUSTR($string$) & adjust right of the string\cr
}
\sms
Here is a program to find out whether a sentence is a palindrome.
\ttbg
PROGRAM PALINDROME
   CHARACTER(LEN = 200) :: TEXT
   INTEGER :: I, J
   LOGICAL :: MATCH

   READ "(A)", TEXT
   J = LEN_TRIM(TEXT)
   MATCH = .TRUE.
   DO I = 1, J/2
      IF(TEXT(I:I) /= TEXT(J:J)) THEN
         MATCH = .FALSE.
         EXIT
      ELSE
         J = J - 1
      END IF
   END DO
   IF (MATCH) THEN
      PRINT *, "Palindrome"
   ELSE
      PRINT *, "Not a palindrome"
   END IF

END PROGRAM
\tted
If {\tt NAT SAW I WAS TAN} is entered, the computer responses with
{\tt Palindrome}.  It says that {\tt MADAM I'M ADAM} is
not a palindrome.

\lecture{8. Pointers}

\section{Pointer declarations, assignments}
Fortran pointers are somewhat different from pointers in other languages
like C or Pascal.  Pointers in Fortran 90 is more conveniently thought
as aliases for other variables.  It is better not to consider pointers
as addresses. To declare pointer variables, we use
\sms
\ll{REAL, POINTER :: P}
\ll{REAL, DIMENSION(:), POINTER :: V}
\ll{INTEGER, POINTER :: IP}
\ll{INTEGER, DIMENSION(:,:), POINTER :: M}
\sms
Then the pointer assignment statements
\sms
\ll{P => R}
\ll{V => A(4, :)}
\ll{IP => IA(1,2)}
\ll{M => IA(1:4, 2:5)}
\sms
associates the pointer {\tt P} with the real variable {\tt R}, or
makes the pointer pointing to the real variable, or maybe more
appropriately stated, {\tt P} is an alias for {\tt R}.  
Here {\tt V} is a pointer of a one-dimensional array and 
is associated with the fourth row of the array {\tt A}.  
Now {\tt M} is an array of $4\times 4$, associated with the array
section {\tt A(1:4, 2:5)}.  {\tt IP} is a single integer variable.
A pointer
of real can not point to any real variables.  In order to make the
above statements legal, we must also declare {\tt R}, {\tt A}, 
and {\tt IA} with attribute {\tt TARGET}:
\sms
\ll{REAL, TARGET :: R}
\ll{REAL, DIMENSION(4, 4), TARGET :: A}
\ll{INTEGER, DIMENSION(10, 5), TARGET :: IA}
\sms
After the pointer assignment, we can use {\tt P} as a real variable and
{\tt V} as a one-dimensional array of four elements.  E.g.,
\sms
\ll{P = 1.0}
\l{\v`PRINT *, V`\m{four numbers will be printed}}
\sms
which is equivalent to 
\sms
\ll{R = 1.0}
\ll{PRINT *, A(4,:)}
\sms
Other examples
\sms
\l{\v`V => A(1:3,2)`\m{{\ninett V} refers to 
{\ninett A(1,2), A(2,2), A(3,2)}}}
\l{\v`M => IA(2:, 3:4)`\m{{\ninett M} refers to the array section}}
Note that in the last example, {\tt M(1,1)} refers to {\tt IA(2,3)},
{\tt M(2,1)} is {\tt IA(3,3)}, and {\tt M(2,2)} for {\tt IA(3,4)}, etc.
An array section or array expression is taken to have lower bounds 1
and upper bounds equal to the extents (dimensions).

\section{Pointer assignment {\tt =>} vs. usual assignment {\tt =}}

The ordinary assignment operator {\tt =} means to copy the value of righthand 
side to the variable on the lefthand side.  The pointer assignment
{\tt =>} means to make the pointer on the lefthand side an alias for
the variable on the righthand side.  E.g.,
\ttbg
REAL, POINTER :: P1, P2
REAL, TARGET :: R1, R2
R1 = 1.1; R2 = 2.2
P1 => R1; P2 => R2
\tted
Then {\tt P1} and {\tt P2} are used as {\tt R1} and {\tt R2},
respectively.  The assignment 
\sms
\ll{P2 = P1}
\sms
copies the value at {\tt R1} to {\tt R2}.  On the other hand, the statement
\sms
\ll{P2 => P1}
\sms
makes {\tt P2} an alias for {\tt P1}, since {\tt P1} is an alias for
{\tt R1}, the above pointer assignment makes both {\tt P1} and {\tt
P2} aliases for {\tt R1}, or pointing to {\tt R1}.  

\section{Allocate and deallocate}
Both scalar variable and array can be dynamically allocated. E.g.,
\ttbg
INTEGER, POINTER :: IP
REAL, DIMENSION(:), POINTER :: V
\tted
\ll{ALLOCATE(IP)}
\ll{ALLOCATE(V(100))}
\sms
or equivalently
\sms
\ll{ALLOCATE(IP, V(100), STAT = $var$)}
\sms
create space for data object that {\tt IP} is pointing to.  And {\tt
IP} is an alias for that space.  And {\tt V} is an array with
100 elements.  {\tt STAT} is an optional keyword argument.
$var$ is zero for successful allocation and positive $var$ indicates
error. It is an error to refer a pointer
without associating the pointer with memory space.   After use
\sms
\ll{DEALLOCATE(IP, V)}
\sms
return the memory space to the system.

\section{Nullify and associated}
In the beginning of the program execution, pointers are undefined.  
To make the pointer a null pointer (pointing nowhere), use the
statement 
\sms
\ll{NULLIFY(P)}
\sms
Later on, the pointer can be tested with the intrinsic function
\sms
\ll{ASSOCIATED(P)}
\sms
which return {\tt .TRUE.} if the pointer is associated with some
object, or {\tt .FALSE.} if it is a null pointer.   The function can
also takes two arguments
\sms
\ll{ASSOCIATED(P, $target$)}
\sms
It returns {\tt .TRUE.} If {\tt P} is associated with the $target$, and
{\tt .FALSE.} otherwise. 



\lecture{9. Structures, Derived Types, and Linked Lists}

Structures and derived types are not available in Fortran 77.
Here is an example of structure definition
\sms
\vbox{\l{\v`TYPE PERSON`\m{the word {\ninett TYPE} plays the r\^ole {\ninett
		   struct} in C}}
\l{\v`   CHARACTER(LEN=40) :: NAME`\m{first field is called {\ninett
			            NAME}, a 40 character string}}
\l{\v`   INTEGER ::ID_NO`\m{second field is {\ninett ID\_NO}, an integer}}
\l{\v`   INTEGER :: AGE`\m{last field is the age}}
\l{\v`END TYPE PERSON`\m{this end the definition}}}
\sms
each of the field can also be an array, or structure, or pointer.  
A field is also called a component in Fortran.
Declarations of the structure type variables are by a statement:
\sms
\ll{TYPE (PERSON) :: STUDENT, TMP}
\sms
which declares that the variables {\tt STUDENT} and {\tt TMP} are of
the type {\tt PERSON}.  We can assign values to the variables with
\sms
\l{\v`STUDENT%NAME = 'Joe Fung'`\m{{\ninett \%} is the structure
				   member operator}}
\l{\v`STUDENT%ID_NO = 7726880`\m{}}
\l{\v`STUDENT%AGE = 19`\m{}}
\sms
The above statements can be replaced by a single statement
\sms
\ll{STUDENT = PERSON('joe Fung', 7726880, 19)}
\sms
Note that the structure name {\tt PERSON} is used as a function that
returns a value of the corresponding type. 
Structure assignment like {\tt TMP = STUDENT} is allowed. It can be
passed to a function as argument and returned as value.  It is possible
to define other operations by operator overloading, such that
{\tt TMP - STUDENT} or {\tt TMP * STUDENT} makes sense.
Here are examples which define point and triangle data types.
\sms
\l{\v`TYPE POINT`\m{define a point as two real numbers}}
\l{\v`   REAL X, Y`\m{}}
\l{\v`END TYPE POINT`\m{}}
\sms
\l{\v`TYPE TRIANLE`\m{a triangle consists of three points}}
\l{\v`   TYPE(POINT) A, B, C`\m{}}
\l{\v`END TYPE TRIANGLE`\m{}}
\sms
We can declare a triangle like this:
\sms
\ll{TYPE(TRIANGLE) :: T}
\sms
{\tt T} has components {\tt T\%A}, {\tt T\%B}, and {\tt T\%C}. All
of type {\tt POINT}, and {\tt T\%A} has components {\tt T\%A\%X},
and {\tt T\%A\%Y} of type {\tt REAL}. Here is some expressions involving
derived types.
\ttbg
T%A = POINT(0.0, 1.0)
T%B%X = 5.0
T%B%Y = T%B%X
\tted

\section{Array of pointers}
In Fortran 90, one can not have array of pointers as such, the following
declaration does not work:
\ttbg
REAL, DIMENSION(10), POINTER :: P
REAL, POINTER :: P(10)
\tted
We can get around using a user-defined type
\ttbg
TYPE ROW
   REAL, POINTER :: R
END TYPE ROW
\tted
\ll{TYPE(ROW), DIMENSION(10) :: P}
\sms
then {\tt P(1)\%R} refers to the first pointer, and 
{\tt P(2)\%R} refers to the second pointer, and so on.

\section{Array as component of derived type}
Look at this definition
\ttbg
TYPE TRIPLET
   REAL :: U
   REAL, DIMENSION(3) :: DU
   REAL, DIMENSION(3,3) :: D2U
END TYPE TRIPLET
\tted
where we can think of {\tt U} as a value of a scalar field, {\tt DU}
is its partial derivative, and {\tt D2U} is the second partial
derivative. With the declaration
\ttbg
TYPE(TRIPLET) :: T
TYPE(TRIPLET), DIMENSION(10,10,10) :: A
\tted
{\tt T\%U} is a real number, {\tt T\%DU} is an array with elements
{\tt T\%DU(1)},  {\tt T\%DU(2)},  and {\tt T\%DU(3)}. {\tt T\%D2U}
is a two-dimensional array with elements {\tt T\%D2U(I,J)}. 
{\tt A(1,2,5)} is of type {\tt TRIPLET}.  {\tt A(1,2,5)\%D2U(1,2)}
refers to element (1,2) of the array {\tt D2U} field in {\tt TRIPLET}
(1,2,5) element of array {\tt A}. 

\section{Linked lists}
A linked list is useful when we don't know the number of elements
ahead, and members of the list may be inserted and removed. 
Each element or node is a structure
\sms
\l{\v`TYPE NODE`\m{}}
\l{\v`   INTEGER :: VALUE`\m{here stores information at each node}}
\l{\v`   TYPE (NODE), POINTER :: NEXT`\m{last item is a pointer to
					 next node}}
\l{\v`END TYPE NODE`\m{}}
\sms
Let's write a subroutine that inserts a node to a list according to 
increasing order of the integer value at the nodes. 
\ttbg
PROGRAM LIST
   IMPLICIT NONE

   TYPE NODE
      INTEGER :: VALUE
      TYPE (NODE), POINTER :: NEXT 
   END TYPE NODE
   TYPE (NODE), POINTER :: L
   INTEGER :: NUMBER, IOS 
 
   ALLOCATE(L)
   NULLIFY(L%NEXT)
   DO
      READ(*, *, IOSTAT = IOS) NUMBER
      IF (IOS < 0) EXIT
      CALL INSERT(L, NUMBER)
   END DO
   CALL PRINT_LIST(L)

CONTAINS

SUBROUTINE INSERT(L, N)
   IMPLICIT NONE
   TYPE (NODE), POINTER :: L, TRAIL, TMP
   INTEGER :: N
   TRAIL => L
   TMP => L%NEXT
   DO
      IF (.NOT. ASSOCIATED(TMP) ) EXIT
      IF (N < TMP%VALUE ) EXIT
      TRAIL => TMP
      TMP => TMP%NEXT
   END DO
   ALLOCATE(TRAIL%NEXT)
   TRAIL%NEXT%VALUE = N
   TRAIL%NEXT%NEXT => TMP
END SUBROUTINE INSERT

SUBROUTINE PRINT_LIST(L)
   IMPLICIT NONE
   TYPE (NODE), POINTER :: L, TMP
   TMP => L%NEXT
   DO
      IF (.NOT. ASSOCIATED(TMP) ) EXIT
      PRINT *, TMP%VALUE
      TMP => TMP%NEXT
   END DO
END SUBROUTINE PRINT_LIST

END PROGRAM LIST
\tted


\lecture{10. Recursion and Trees}

\section{Recursive functions and subroutines}
Normally, a function or subroutine can not call itself recursively.  
This is the case in Fortran~77. The new Fortran~90 allows recursive
calls, but the function or subroutine must be prefixed with the keyword
{\tt RECURSIVE}.  In addition, the recursive function definition must
specify a different name for the result, e.g.,
\ttbg
RECURSIVE FUNCTION FACTORIAL(N) RESULT (FAC)
   INTEGER, INTENT(IN) :: N
   INTEGER :: FAC
   IF (N <= 0) THEN
      FAC = 1
   ELSE
      FAC = N * FACTORIAL(N - 1)
   END IF
END FUNCTION FACTORIAL
\tted
Note that the return value of the function is assigned to the result
{\tt FAC}, and {\tt FACTORIAL} should not be used for this purpose.
The attribute {\tt INTENT(IN)} indicates that {\tt N} should not be
modified in the function.  

The structure of a recursive subroutine is like this
\ttbg
RECURSIVE SUBROUTINE SUB( ... )
   ...
   CALL SUB( ... )
   ...
END SUBROUTINE SUB
\tted

\section{Tree}
A basic node of a binary tree can be declared as
\sms
\l{\v`TYPE NODE`\m{name of the structure}}
\l{\v`   INTEGER :: VALUE`\m{value stored at each node}}
\l{\v`   TYPE (NODE), POINTER :: LEFT, RIGHT`\m{pointers to children}}
\l{\v`END TYPE NODE`\m{}}
\sms

The following program involves two recursive subroutines, {\tt INSERT} 
subroutine inserts a number such that at each node the number
at left child is less than that of the node, and the number at right
child is equal to or greater than the number at the node. The {\tt
PRINT\_TREE} function prints out the number {\sl inorder} -- that is, go to
left, print the current node, go to right.
\ttbg
PROGRAM TREE
   IMPLICIT NONE

   TYPE NODE
      INTEGER :: VALUE
      TYPE (NODE), POINTER :: LEFT, RIGHT
   END TYPE NODE

   TYPE (NODE), POINTER :: T
   INTEGER :: NUMBER, IOS
   
   NULLIFY(T)
   DO 
      READ(*, *, IOSTAT = IOS) NUMBER
      IF ( IOS < 0 ) EXIT
      CALL INSERT(T, NUMBER)
   END DO

   CALL PRINT_TREE(T)

CONTAINS

RECURSIVE SUBROUTINE INSERT(T, N)
   IMPLICIT NONE
   TYPE (NODE), POINTER :: T
   INTEGER :: N
   IF (.NOT. ASSOCIATED(T)) THEN
      ALLOCATE(T)
      T%VALUE = N
      NULLIFY(T%LEFT)
      NULLIFY(T%RIGHT)
   ELSE IF (N < T%VALUE) THEN
      CALL INSERT(T%LEFT, N)
   ELSE
      CALL INSERT(T%RIGHT, N)
   END IF
END SUBROUTINE INSERT

RECURSIVE SUBROUTINE PRINT_TREE(T)
   IMPLICIT NONE
   TYPE (NODE), POINTER :: T
   IF (ASSOCIATED(T)) THEN 
      CALL PRINT_TREE(T%LEFT)
      PRINT *, T%VALUE
      CALL PRINT_TREE(T%RIGHT)
   END IF
END SUBROUTINE PRINT_TREE

END PROGRAM TREE
\tted

\lecture{11. Modules}

\section{data sharing}

The concept of module is new.  Fortran~77 does not support modules.
It has a number of uses, one of them is to share variables between
different subprograms.  There are four ways of passing data between
subprogram units.  

(1) Through arguments of function and subroutine, both for passing in and out.

(2) Use the data in host program directly by internal subprograms.
This is not possible in Fortran~77.

(3) Use common blocks between subprograms.  Beside using arguments, this 
is the only option available in Fortran~77.  And it is considered obsolete in
Fortran~90. 

(4) Use module provided by Fortran~90.

Let's look at each case in detail.  When using arguments,
we must remember that all arguments are passed ``call-by-reference''.
The effect is that variables modified by the subprogram are
passed back to the calling program.  

Here is an example that uses global variables.
\sms
\l{\v`PROGRAM`\m{}}
\l{\v`   INTEGER :: I, J, K`\m{declare global variables}}
\l{\v`   I = 1; J = 10; K = 3`\m{}}
\l{\v`   ...`\m{}}
\l{\v`   CALL SUB(Y)`\m{}}
\l{\v`   ...`\m{}}
\l{\v`CONTAINS`\m{internal subprograms start here}}
\l{\v`   SUBROUTINE SUB(X)`\m{}}
\l{\v`      X = I + J`\m{global variables are used here}}
\l{\v`      K = K + 1`\m{}}
\l{\v`      ...`\m{}}
\l{\v`   END SUBROUTINE SUB`\m{}}
\l{\v`   ...`\m{}}
\l{\v`END PROGRAM`\m{}}
\sms

We now explain the common blocks. Suppose one of the
subroutine wants to have some variables shared by other subroutines
or functions, after the declaration, one can put
\sms
\ll{COMMON/$cmname$/ SIZE, L, A(10)}
\sms
where $cmname$ delimited by slashes is a unique global identifier to 
specify the common block, following it are the names of variables.
In other subroutines, one can put the same line of common block
specification, then the variables {\tt SIZE}, {\tt L}, and {\tt A} 
in different subprograms refer to the same memory locations.  They
share the {\sl common} data.  The order of the variables 
in common block is important,
but you can have different names in different subprograms for the
variables. 

A better method is provided in Fortran 90 with the modules.  E.g.,
\ttbg
MODULE COMMON_DATA
   INTEGER :: SIZE, L
   REAL, DIMENSION (10) :: A
END MODULE COMMON_DATA
\tted
A {\tt USE} statement will make those data available in a subprogram
without declaration, e.g.,
\ttbg
SUBROUTINE TEST_USE
   USE COMMON_DATA
   A(1) = 1.0
   ...
\tted
\sms

\section{Procedures contained in module}
A module can contain function and subroutine definitions as well as
data type definitions, declarations, and interface specifications.
These definitions then can be used by other program units.  The
general form of a module is 
\sms
\ll{MODULE $module$\_$name$}
\ll{\ \ \ $type$ $declarations,$ $etc.$}
\ll{\ \ \ CONTAINS}
\ll{\ \ \ $function$ $and$ $subroutine$ $definitions$}
\ll{END MODULE $module$\_$name$}
\sms

\section{Public and private attributes}
By default, the variables declared in a module are public, so that
the program units  using the module can also use these variables.
Private variables are local to the module.  The module functions
contained in the module
can use these variables, but they are invisible to program units
which use the module.  Here are three different ways of declaring
the public and private attributes.
\ttbg
REAL, PUBLIC :: X, Y, Z
INTEGER, PRIVATE :: U, V, W
\tted
or
\ttbg
REAL :: X, Y, Z
INTEGER :: U, V, W
PUBLIC :: X, Y, Z
PRIVATE :: U, V, W
\tted
or
\ttbg
PUBLIC
   REAL :: X, Y, Z
PRIVATE
   INTEGER :: U, V, W
\tted
Besides the public and private module variables, the module
functions and subroutines can also be declared as public or private.

\section{The use statement}
The statement
\sms
\ll{USE $module$\_$name$}
\sms
provide access to all the data objects, derived types, interface
blocks, procedures, etc.  {\tt USE} statement must be the first
statement in a program unit.  E.g.,
\ttbg
USE COMMON_DATA
\tted
The names in module may be renamed
\sms
\l{\v`USE COMMON_DATA, B => A`\m{variable {\ninett A} is renamed as {\ninett B}}}
\sms
Or only a selected names are accessed
\sms
\l{\v`USE COMMON_DATA, ONLY: L, A`\m{other data items can not be accessed}}
\sms

\section{Operator overloading and user-defined operators}
In Fortran 90, it is possible to extend the meaning of the 
intrinsic operators (such as {\tt +}, {\tt -}, {\tt *}, {\tt /}, {\tt **}, 
{\tt //}, {\tt .EQ.}) by defining an interface block together with a 
function which defines the operation inside a module.  For example, 
we like to use asterisk for  matrix multiplication, we first specify an 
interface
\ttbg
INTERFACE OPERATOR (*)
   MODULE PROCEDURE PROD
END INTERFACE
\tted
where {\tt PROD} is a function which has two arguments, we'll define
this function later.  Once our module is ready, we can use it.
For example, we could specify a matrix multiplication as
\sms
\ll{C = A * B}
\sms
where {\tt C} is a matrix of dimension $m \times n$, {\tt A} and {\tt
B} are $m\times k$ and $k \times n$, respectively. 

It is also possible to define an operator of your own, e.g., an
inverse operator for matrix.  Use-defined operator is  a
name delimited by periods, like {\tt .INVERSE.}, and is specified with 
\ttbg
INTERFACE OPERATOR (.INVERSE.)
   MODULE PROCEDURE MINVERSE
END INTERFACE
\tted
The function {\tt MINVERSE} which takes only one argument could be very long. 
To invert a matrix {\tt A} and assign it to {\tt C} we simply write
\sms
\ll{C = .INVERSE. A}
\sms

\section{Generic assignment}
One can also extend the meaning of an assignment with an interface block
\sms
\ll{INTERFACE ASSIGNMENT (=)}
\ll{\ \ \ MODULE PROCEDURE $subroutine$\_$name$}
\ll{END INTERFACE}
\sms
where the subroutine defined in the same module takes two arguments,
the first is the left-hand side of the equal sign, and the second is
the right-hand side of the equal sign. 

In the operator overloading and assignment specification, the {\tt 
MODULE PROCEDURE} keywords should be replaced by function or
subroutine specification like the usual interface specification if
the function or subroutine are external.  For example, if the {\tt PROD}
function is written as an external function, then we need the following
interface specification to overload the multiplication operator {\tt *},
\ttbg
INTERFACE OPERATOR (*)
   FUNCTION PROD(A, B)
      TYPE(MATRIX) :: A, B, PROD
   END FUNCTION PROD
END INTERFACE
\tted

\section{Summary of interface specifications}
(1) Interface declares the arguments and function return value types, like
the C prototype declarations. E.g.,
\ttbg
INTERFACE
   FUNCTION MINCON(N, X, UPPER, LOWER)
      INTEGER :: N, UPPER
      REAL :: MINCON, X
      REAL, OPTIONAL :: LOWER
   END FUNCTION MINCON
END INTERFACE
\tted

(2) Specify generic names
\ttbg
INTERFACE GAMMA
   FUNCTION SGAMMA(X)
      REAL(SELECTED_REAL_KIND(6)) :: SGAMMA, X
   END
   FUNCTION DGAMMA(X)
      REAL (SELECTED_REAL_KIND(12)) :: DGAMMA, X
   END
END INTERFACE
\tted
With this interface specification, when {\tt GAMMA} is used in the 
program, either {\tt SGAMMA} (a single precision version) or
{\tt DGAMMA} (a double precision version) will be used depending on 
the data type of the argument {\tt X} given.

(3) Operator overloading
\sms
\l{\v`INTERFACE OPERATOR (+)`\m{require a two-argument function}}
\ll{INTERFACE OPERATOR (.$op$.)\m{one or two arguments function}}
\l{\v`INTERFACE ASSIGNMENT (=)`\m{two arguments subroutine}}

\section{An example of a module}

Below is a complete example which shows various aspects of a module:
data type definitions, interfaces, module procedures.  The module
{\tt MATRICES} defines several important operations on matrices.
\sms
\l{\v`MODULE MATRICES`\m{module program unit start}}
\l{\v` `\m{}}
\l{\v`   TYPE MATRIX`\m{matrix data type is defined as}}
\l{\v`      INTEGER :: ROW`\m{a structure containing the number of rows}}
\l{\v`      INTEGER :: COL`\m{and columns, as integer variables,}}
\l{\v`      REAL, DIMENSION(:, :), POINTER :: ARRAY`\m{and a pointer to a
						two-dimensional array}}
\l{\v`   END TYPE MATRIX`\m{}}
\l{\v``\m{}}
\l{\v`   INTERFACE OPERATOR (+)`\m{the function {\ninett SUM} can be used}}
\l{\v`      MODULE PROCEDURE SUM`\m{through the addition operator {\ninett +}}}
\l{\v`   END INTERFACE`\m{by an interface specification}}
\l{\v``\m{}}
\l{\v`   INTERFACE OPERATOR (*)`\m{asterisk means both matrix multiplication}}
\l{\v`      MODULE PROCEDURE PROD, SCALE`\m{and multiplying matrix with
					    a number}}
\l{\v`   END INTERFACE`\m{}}
\l{\v``\m{}}
\l{\v`   PRIVATE SUM, PROD, SCALE`\m{make {\ninett SUM}, 
{\ninett PROD}, and {\ninett SCALE} local to this module}}
\l{\v` `\m{}}
\l{\v`CONTAINS`\m{the module procedures follow}}
\l{\v` `\m{}}
\l{\v`   FUNCTION SUM(A, B) `\m{}}
\l{\v`      TYPE(MATRIX) :: A, B, SUM`\m{}}
\l{\v`      IF (A%ROW /= B%ROW .OR. A%COL /= B%COL ) &`\m{}}
\l{\v`         STOP 'Can not add'`\m{abort when error occurs}}
\l{\v`      SUM%ROW = A%ROW`\m{}}
\l{\v`      SUM%COL = A%COl`\m{}}
\l{\v`      ALLOCATE( SUM%ARRAY(A%ROW, A%COL) )`\m{}}
\l{\v`      SUM%ARRAY = A%ARRAY + B%ARRAY`\m{}}
\l{\v`   END FUNCTION SUM`\m{}}
\l{\v` `\m{}}
\l{\v`   FUNCTION PROD(A, B)`\m{}}
\l{\v`      TYPE(MATRIX) :: A, B, PROD`\m{}}
\l{\v`      IF (A%COL /= B%ROW ) STOP 'Can not multiply'`\m{}}
\l{\v`      PROD%ROW = A%ROW`\m{}}
\l{\v`      PROD%COL = B%COL`\m{}}
\l{\v`      ALLOCATE( PROD%ARRAY(A%ROW, B%COL) )`\m{}}
\l{\v`      PROD%ARRAY = MATMUL(A%ARRAY, B%ARRAY)`\m{{\ninett MATMUL} is 
					an intrinsic function}}
\l{\v`   END FUNCTION PROD`\m{}}
\l{\v`   `\m{}}
\l{\v`   FUNCTION SCALE(C, A)`\m{multiply a number to a matrix}}
\l{\v`      TYPE(MATRIX) :: A, SCALE`\m{}}
\l{\v`      REAL :: C`\m{}}
\l{\v`      SCALE%ROW = A%ROW`\m{}}
\l{\v`      SCALE%COL = A%COL`\m{}}
\l{\v`      ALLOCATE( SCALE%ARRAY(A%ROW, A%COL) )`\m{}}
\l{\v`      SCALE%ARRAY = C * A%ARRAY`\m{}}
\l{\v`   END FUNCTION SCALE`\m{}}
\l{\v` `\m{}}
\l{\v`   SUBROUTINE SETM(ROW, COLUMN, ARRAY, A)`\m{initialize a matrix}}
\l{\v`      TYPE(MATRIX) :: A`\m{}}
\l{\v`      INTEGER :: ROW, COLUMN`\m{}}
\l{\v`      REAL, DIMENSION(:,:) :: ARRAY`\m{}}
\l{\v`      A%ROW = ROW`\m{}}
\l{\v`      A%COL = COLUMN`\m{}}
\l{\v`      ALLOCATE( A%ARRAY(A%ROW, A%COL) )`\m{}}
\l{\v`      A%ARRAY = ARRAY`\m{}}
\l{\v`   END SUBROUTINE SETM`\m{}}
\l{\v` `\m{}}
\l{\v`   SUBROUTINE PRINTM(A)`\m{print a matrix}}
\l{\v`      TYPE (MATRIX) :: A`\m{}}
\l{\v`      INTEGER :: I, J`\m{}}
\l{\v`      DO I = 1, A%ROW`\m{}}
\l{\v`         PRINT *, ( A%ARRAY(I,J), J = 1, A%COL )`\m{}}
\l{\v`      END DO`\m{}}
\l{\v`   END SUBROUTINE PRINTM`\m{}}
\l{\v` `\m{}}
\l{\v`END MODULE MATRICES`\m{}}
\sms 
You might want to add other operations like transpose, matrix
inverse, or determinant to  the module.   The use of memory is not efficient
in these functions.  We asked for memory to hold every immediate result,
but we couldn't return the memory to the system. Here is a test program that
uses the module. 
\sms
\l{\v`PROGRAM TEST`\m{}}
\sms
\l{\v`USE MATRICES`\m{}}
\sms
\l{\v`  TYPE(MATRIX) :: A, B, C`\m{}}
\l{\v`  REAL, DIMENSION(2,4) :: TMP`\m{}}
\l{\v`  TMP(1,:) = (/1.0, 2.0, 3.0, 4.0 /)  `\m{}}
\l{\v`  TMP(2,:) = (/5.0, 6.0, 7.0, 8.0 /)`\m{}}
\l{\v`  CALL SETM(2, 2, TMP(1:2,1:2), A)`\m{}}
\l{\v`  CALL SETM(2, 4, TMP, B)`\m{}}
\l{\v`  PRINT *, 'matrix A'`\m{}}
\l{\v`  CALL PRINTM(A)`\m{}}
\l{\v`  PRINT *, 'matrix B'`\m{}}
\l{\v`  CALL PRINTM(B)`\m{}}
\l{\v`  C = A*B + 2.0 * B`\m{}}
\l{\v`  PRINT *, 'matrix AB + 2B'`\m{}}
\l{\v`  CALL PRINTM(C)`\m{}}
\sms
\l{\v`END PROGRAM TEST`\m{}}
\sms

\section{Attributes vs. statements}
In general, variable declaration takes the form
\sms
\ll{$type$, $attribute$ :: $variable$}
\sms
For example
\sms
\l{\v`REAL, PARAMETER :: PI = 3.141592653`\m{{\ninett PARAMETER} attribute}}
or equivalently\par
\ll{REAL :: PI}
\l{\v`PARAMETER (PI = 3.141592653)`\m{{\ninett PARAMETER} statement}}
You may specify with attributes
\ttbg
REAL, ALLOCATABLE :: A(:,:)
REAL, POINTER :: SON
REAL, TARGET :: Y(10)
REAL, SAVE :: S
\tted
or equivalently, use corresponding statements
\ttbg
REAL :: A, SON, Y, S
ALLOCATABLE :: A(:,:)
POINTER :: SON
TARGET :: Y(10)
SAVE :: S
\tted
{\tt SAVE} attribute is like static variable in C function.  The content
will be preserved between function calls.

\section{DATA statement}
{\tt DATA} statement is considered obsolete in Fortran 90.  It is used
in Fortran 77 to initialize data.  E.g.,
\ttbg
REAL :: A = 1.0, B = 2.0, C = 3.0
INTEGER :: I = 1, J = 2, K = 3
\tted
or equivalently
\ttbg
REAL :: A, B, C
INTEGER :: I, J, K
DATA A, B, C /1.0, 2.0, 3.0/, I, J, K /1, 2, 3/
\tted
other {\tt DATA} statements
\sms
\l{\v`DATA I/0/`\m{}}
\l{\v`DATA I, J, K /3*0/`\m{all three variables initialize to 0}}
\l{\v`DATA R(1:LENGTH) /LENGTH*0./`\m{all elements gets 0.0}}
\sms
Array elements or array sections may initialized this way, eg.,
\sms
\l{\v`REAL A(5,5)`\m{a 5 by 5 matrix}}
\l{\v`DATA A(1,1) /1.0/`\m{a single element}}
\l{\v`DATA A(2:5, 4) /4*1.0/`\m{array section with 4 elements}}
\l{\v`DATA A /25*1.0/`\m{all array elements are set to 1.0}}
\sms
Implicit loops may be used
\sms
\ll{DATA (( A(I,J), I = 1, 5, 2), J = 1, 5) /15*0./}
\sms


\lecture{12. Intrinsic Procedures}
Fortran has a large collections of so-called intrinsic functions and
subroutines for scientific computation.  This is one of the advantage of
using Fortran for scientific and engineering calculations. 

The function names are generic names.  This means that the function works
for integer, real, double, or complex argument.
Most of the functions also work on arrays.  E.g.,
\ttbg
X = (/ 1.0, 4.0, 9.0, 16.0 /)
Y = SQRT(X)
\tted
then {\tt Y} has value {\tt (/1.0, 2.0, 3.0, 4.0/)}. 
There are two more statements to do with intrinsic functions.
\sms
\ll{INTRINSIC $intrinsic$\_$procedure$\_$name$}
\sms
inform the compiler (as well as the programmer) the name refers to
intrinsic function or subroutine.
\sms
\ll{EXTERNAL $name$}
\sms
Specify the name is an external function or subroutine.

\section{Mathematical and numerical functions}
\sms
\halign{&\t\tt#\hfil&\ $#$\hfil&\ #\hfil\cr
ACOS(X) & \cos^{-1} x & arccosine & ASIN(X) & \sin^{-1} x & arcsine \cr
COS(X) & \cos x & cosine & SIN(X) & \sin x & sine \cr
ATAN(X) & \tan^{-1} x & arctangent & TAN(X) & \tan x & tangent \cr
TANH(X) & \tanh x & hyperbolic tangent & LOG10(X) & \log_{10} x & common logarithm \cr
EXP(X) & e^{x} & exponential & LOG(X) & \ln x & natural logarithm \cr
COSH(X) & \cosh x & hyperbolic cosine & SINH(X) & \sinh x & hyperbolic sine \cr
SQRT(X) & \sqrt{x} & square root & ABS(X) & |x| & absolute value \cr
CEILING(X) & \lceil x \rceil & ceiling & FLOOR(X) & \lfloor x \rfloor & floor \cr
MAX(A1,A2,...) & & maximum & MIN(A1,A2,...) & & minimum \cr
INT(X) & & convert to integer & NINT(X) & & nearest integer \cr
MOD(A,P) & a {\rm\ mod\ } p & remainder & SIGN(A,B) & & transfer of sign \cr
}
\sms

\section{Bit manipulation functions}
\sms
\halign{&\t\tt#\hfil&\qquad #\hfil\cr
BTEST(I, POS) & bit testing & IBCLR(I, POS) & clear bit \cr
IBSET(I, POS) & set bit & IBITS(I, POS, LEN) & bit extract \cr
IAND(I, J) & bitwise logical and & IOR(I, J) & bitwise or \cr
IEOR(I, J) & bitwise exclusive or & NOT(I) & complement \cr
ISHFT(I, SHIFT) & bit shift & ISHFTC(I, SHIFT) & circular shift \cr
}
\sms

\section{Array functions}

Fortran 90 has a large number of intrinsic functions on array.
\sms
\halign{&\t\tt#\hfil&\  #\hfil\cr
DOT\_PRODUCT(A, B) & dot product of vectors & MATMUL(A, B) & matrix multiplication \cr
MAXVAL(A) & maximum value in array & MINVAL(A) & minimum value in array \cr
MAXLOC(A) & location of maximum & MINLOC(A) & location of minimum \cr
PRODUCT(A) & product of array elements & SUM(A) & sum of array elements \cr
TRANSPOSE(A) & transpose of matrix &  CSHIFT(A,SHIFT,DIM) & circular shift \cr
SIZE(A) & total number of elements & SHAPE(A) & shape of array \cr
LBOUND(A) & lower dimension bounds & UBOUND(A) & upper dimension bounds \cr
RESHAPE(SOURCE,SHAPE) & reshape array & COUNT(MASK) & number of true elements \cr 
}
\sms
Example:
\ttbg
REAL, DIMENSION(-10:5, 2) :: A
REAL R
\tted
\l{\v`SIZE(A)`\m{has value 32}}
\l{\v`SIZE(A, 1)`\m{has value 16}}
\l{\v`SIZE(A, 2)`\m{has value 2}}
\l{\v`LBOUND(A)`\m{has value {\ninett (/-10, 1/)}}}
\l{\v`UBOUND(A)`\m{has value {\ninett (/5, 2/)}}}
\l{\v`SHAPE(A)`\m{has value {\ninett (/16, 2/)}}}
\l{\v`SHAPE(R)`\m{has value 0}}
\sms

\section{Numeric inquiry functions}
\halign{\t\tt#\hfil&\qquad \tt# \hfil&\qquad #\hfil\cr
$function(arg)$ & $typical$\ $value$ & $explanation$ \cr
\noalign{\sms}
DIGITS(1) & 31 & binary digits for an integer\cr
DIGITS(1.0) & 24 & or real representation \cr
EPSILON(1.0) & 1.2E-7 & number negligible compares to 1.0\cr
HUGE(1) & 2147483647 & largest number of given type \cr
HUGE(0.0) & 3.4E38 & \cr
MAXEXPONENT(1.0) & 127 & maximum and minimum \cr
MINEXPOMENT(1.0) & -126 & exponent value (base 2) \cr
PRECISION(1.0) & 6 & precision of real \cr
RADIX(1.0) & 2 & base of number \cr
RADIX(1) & 2 & representation \cr
RANGE(1) & 9 & range of number $-10^9$ to $10^9$ \cr
RANGE(1.0) & 37 & for real $10^{-37}$ to $10^{37}$ \cr
TINY(1.0) & 1.1E-38 & smallest positive number \cr
}

\section{Other intrinsic subroutines}
\sms
\halign{\t\tt#\hfil&\qquad #\hfil\cr
DATE\_AND\_TIME(ALL, COUNT, MSECOND, SECOND, & date and time \cr
\ \ \ \ MINUTE, HOUR, DAY, MONTH, YEAR, ZONE)  & all arguments are optional \cr
RANDOM(R) & return random number $0 < \hbox{\tt R} <1$ \cr
RANDOM\_SEED() & initialize or restart random number generator \cr
SYSTEM\_CLOCK(COUNT,COUNT\_RATE,COUNT\_MAX) & data from system clock \cr
}


\vfill\eject
\lecture{13. Introduction to MPI}

\section{Parallel Computation}

Parallel computation is playing more and more important rule in
large-scale scientific and engineering computing.   The speed of
a single processor is fundamentally limited.  Using many processors,
say from few upto a thousand, to solve a single problem is possible 
now a days. 

Two distinct computer architectures exit: distributed memory and
shared memory.  The acronym  MPI is from ``Massage-Passing Interface''. 
MPI is a software standard and a library to facilitate
parallel computing in distributed memory environment.   With MPI
or other software (e.g., PVM), we can use a workstation cluster
to do parallel computation. 


\section{Computing Pi}

Consider the numerical integration of   
$$ 4 \int_0^1 { 1 \over 1 + x^2 } dx. $$  
The result is $\pi \approx 3.14159$.  The integral is approximated
by simple discrete sum  $  \sum 4/(1+x^2) \Delta x$.   The job of
computing the sum is divided among processors.  Here is the code.
We'll explain them line by line.

\ttbg
c**********************************************************************
c   pi.f - compute pi by integrating f(x) = 4/(1 + x**2)     
c     
c   Each node: 
c    1) receives the number of rectangles used in the approximation.
c    2) calculates the areas of it's rectangles.
c    3) Synchronizes for a global summation.
c   Node 0 prints the result.
c
c  Variables:
c
c    pi  the calculated result
c    n   number of points of integration.  
c    x           midpoint of each rectangle's interval
c    f           function to integrate
c    sum,pi      area of rectangles
c    tmp         temporary scratch space for global summation
c    i           do loop index
c***********************************************************************
      program main

      include 'mpif.h'

      double precision  PI25DT
      parameter        (PI25DT = 3.141592653589793238462643d0)

      double precision  mypi, pi, h, sum, x, f, a
      integer n, myid, numprocs, i, rc
c                                 function to integrate
      f(a) = 4.d0 / (1.d0 + a*a)

      call MPI_INIT( ierr )
      call MPI_COMM_RANK( MPI_COMM_WORLD, myid, ierr )
      call MPI_COMM_SIZE( MPI_COMM_WORLD, numprocs, ierr )
      print *, "Process ", myid, " of ", numprocs, " is alive"

      sizetype   = 1
      sumtype    = 2
      
 10   if ( myid .eq. 0 ) then
         write(6,98)
 98      format('Enter the number of intervals: (0 quits)')
         read(5,99) n
 99      format(i10)
      endif
      
      call MPI_BCAST(n,1,MPI_INTEGER,0,MPI_COMM_WORLD,ierr)

c                                 check for quit signal
      if ( n .le. 0 ) goto 30

c                                 calculate the interval size
      h = 1.0d0/n

      sum  = 0.0d0
      do 20 i = myid+1, n, numprocs
         x = h * (dble(i) - 0.5d0)
         sum = sum + f(x)
 20   continue
      mypi = h * sum

c                                 collect all the partial sums
      call MPI_REDUCE(mypi,pi,1,MPI_DOUBLE_PRECISION,MPI_SUM,0,
     $     MPI_COMM_WORLD,ierr)

c                                 node 0 prints the answer.
      if (myid .eq. 0) then
         write(6, 97) pi, abs(pi - PI25DT)
 97      format('  pi is approximately: ', F18.16,
     +          '  Error is: ', F18.16)
      endif

      goto 10

 30   call MPI_FINALIZE(rc)
      stop
      end
\tted

The statement
\sms
\ll{include 'mpif.h'}
\sms
is necessary in every MPI Fortran program and subprogram to define
various constants and variables.   After the variable definitions,
the three lines appears in every MPI program:
\ttbg
      call MPI_INIT( ierr )
      call MPI_COMM_RANK( MPI_COMM_WORLD, myid, ierr )
      call MPI_COMM_SIZE( MPI_COMM_WORLD, numprocs, ierr )
\tted
The call to {\tt MPI\_INIT} is required and must be the first MPI
call. It establishes the MPI ``environment.''   The last argument
of every MPI subroutine is the returned error code.  The value will
be {\tt MPI\_SUCCESS} if execution is successful.  

All MPI communication is associated with a {\sl communicator} which
defines a group of processes.  {\tt MPI\_COMM\_WORLD} defines the 
set of all processes available.  The return value {\tt myid} from
calling {\tt MPI\_COMM\_RANK} identify the process; {\tt myid} is 
a number from 0 to $1 - {\tt numprocs}$, where {\tt numprocs}
is the total number of processors.  

The processor 0 is normally directly connected with user, and is called
the master process.  The master process gets a value for $n$, the
number of integration points from the user.  The line
\ttbg
      call MPI_BCAST(n,1,MPI_INTEGER,0,MPI_COMM_WORLD,ierr)
\tted
sends the value $n$ to all other processes.   The {\tt MPI\_BCAST} results
in every process ending up with a copy of $n$.    The data to be
communicated is described by the address ({\tt n}), the number of items 
(1), and the datatype ({\tt MPI\_INTEGER}).  The process with the original
copy is specified by the fourth argument (0 in this case, the
master process, which just reads it from the user).   The master 
process 0 sends the data to all other processes.  All other processes
receive the data.  

After the call to {\tt MPI\_BCAST}, all processes have $n$ and their
own identifiers, which is enough information for each one to compute
its contribution, {\tt mypi}.   Each process computes the area of
every {\tt numprocs}s'th rectangle, starting with ${\tt myid} + 1$. 
Next, all of the values of {\tt mypi} held by the individual processes
need to be added up.  MPI provides a rich set of such operations, using
the {\tt MPI\_REDUCE} routine, with an argument specifying which arithmetic
or logical operation is being requested.  In our case the call is
\ttbg
      call MPI_REDUCE(mypi,pi,1,MPI_DOUBLE_PRECISION,MPI_SUM,0,
     $     MPI_COMM_WORLD,ierr)
\tted
The first two arguments identify the source and result addresses,
respectively.  The data being collected consists of 1 item of type
{\tt MPI\_DOUBLE\_PRECISION}.   The operation is addition ({\tt MPI\_SUM}),
and the result of the operation is to be placed in {\tt pi} in the 
process with rank 0 (fifth argument).  The last two arguments  are
the communicator and error code, as usual. 

All processes then return to the top of the loop (the master prints
the answer first).   The {\tt MPI\_BCAST} causes all the processes
except the master to wait for the next value of $n$.   When the user
types a zero in response to the request for a number of rectangles,
the loop terminates and all processes execute
\ttbg
      call MPI_FINALIZE(rc)
\tted
This call must be made by every process in an MPI computation.  It 
terminates the MPI ``environment''; no MPI calls may be made by a 
process after its call to {\tt MPI\_FINALIZE}.  

\section{Running MPI programs}

To run the multi-processor program, we must need to compile the
program with 
\sms
\ll{mpif77 $file$.f}
\sms
\noindent and make sure you have a {\tt .rhosts} file with these machines
as trusted hosts (see manpage on rhosts for more information).   This file
allow you to rsh or rlogin without a password.  Then one could run the
program as
\sms
\ll{mpirun -np 4 $file$}
\sms
where {\tt -np 4} specifies that you use 4 processors,
$file$ is the name of executable.  If the PVM's bin directory 
is not on your search path, you may either add the directory on your
search path or using the full path names in the above examples.

\section{A `Hello World' program in MPI}
The program below shows the fundamentals of a message-passing
parallel program, sending and receiving information. 

\ttbg
C     Each process sends a message back to the master process 0
C     the master prints out the message

      program main

      include 'mpif.h'

      integer revid, myid, numprocs, i
      integer rc, tag, src, ierr
      integer status(MPI_STATUS_SIZE)

      call MPI_INIT( ierr )
      call MPI_COMM_RANK( MPI_COMM_WORLD, myid, ierr )
      call MPI_COMM_SIZE( MPI_COMM_WORLD, numprocs, ierr )

      if (myid .GT. 0) then
         tag = myid 
         call MPI_SEND( myid, 1, MPI_INTEGER, 0, tag, 
     &                 MPI_COMM_WORLD, ierr )
      else
         do 10 i = 1, numprocs - 1
            src = MPI_ANY_SOURCE
            tag = MPI_ANY_TAG
            call MPI_RECV( revid, 1, MPI_INTEGER, src, tag, 
     &                    MPI_COMM_WORLD, status, ierr )
            print *, "Received from process ", revid
10       continue
      endif

      call MPI_FINALIZE(rc)
      stop
      end
\tted
The {\tt MPI\_SEND} and {\tt MPI\_RECV} are the most basic message 
passing functions.   The three arguments of send and receive describe
the message by address, count, and datatype.  The fourth argument 0 in 
{\tt MPI\_SEND} is the destination.  The next argument is an integer
message type, or tag.  We use the tag to indicate extra information
about the message, it is otherwise arbitrary.  {\tt MPI\_COMM\_WORLD} 
is the default communicator. The last one is error return code. 

The receiving function {\tt MPI\_RECV} has a source (fourth) and
tag (fifth argument).  The source indicates from which process
one wants to receive, and the tag indicates from which tag to receive. 
The receiving function will wait until the expected data is 
received.  In the above example, we are not selective, and accept
message from any source or tag.

For complete information on MPI, the reader is recommended to read
the text book by Gropp et al, and the user manual of MPI.

\vfill\eject

\font\bigsc=cmcsc10 scaled\magstep1
\baselineskip = 11pt
\def\]{\leavevmode\hbox{\tt\char`\ }}
\def\fc#1#2{\hskip 3 mm {\tt #1} \hfill $#2$\hskip 3 mm}
\def\fct#1#2{\hskip 3 mm {\tt #1} \hfill #2\hskip 3 mm}
\def\fcc#1#2{\line{\hskip 3 mm {\tt #1} \hfill #2\hskip 3 mm }}
\def\sect#1{\smallskip{\bf #1}\par\nobreak}
\def\tab{\phantom{\]\]\]\]\]\]}}
\def\logo{$^{\scriptscriptstyle 90}$}
\def\shrink{\vskip -2pt}
\parindent=0pt

\centerline{\bigsc Fortran Language Brief Summary}
\medskip
\sect{Compile and Run Fortran Programs}
\fcc{f77 $filename$.f\ \ \ {\rm or} \ \ \ f90 $file$.f90}{compile and link, produce {\tt a.out}}
\fcc{a.out}{execute the program}

\sect{Program Structure}

\fcc{\]\]\]\]\]\]PROGRAM {\it Name}}{specify the main program name (optional)}
\fcc{C\ \ \ \ \ This is a comment}{a line starting with {\tt C} in column 1
is a comment line}
\fcc{\tab {\it variable declarations}}{declare variable types and array types}
\fcc{\tab {\it statements}}{statements are in column 7 to 72}
\fcc{12345 }{columns 1 to 5 are for statement label}
\fcc{\ \ \ \ \ +}{character in column 6 indicates line continuation}
\fcc{\tab READ *, Y}{read input from terminal}
\fcc{\tab CALL SUB(A,B)}{subroutine call}
\fcc{\tab X = FUNC(Y)}{function call}
\fcc{\tab ...}{and so on}
\fcc{\tab PRINT *, ' X = ', X}{print out value, characters are quoted
				with {\tt '...'}}
\fcc{\tab END}{terminate main program}
\fcc{\tab}{}

\fcc{\tab FUNCTION FUNC(X)}{define a function}
\fcc{\tab ... }{ }
\fcc{\tab FUNC = A * X**2 + B/X - C}{
assign value of $A X^2 + B/X - C $ to $FUNC$}
\fcc{\tab END}{end of the function definition}
\fcc{\tab}{}

\fcc{\tab SUBROUTINE SUB(A,B)}{define a subroutine}
\fcc{\tab ... }{variables in subroutine are local}
\fcc{\tab END}{end of the definition}
\fcc{\tab }{ }

\sect{Variable Declarations}
\fcc{\tab INTEGER I,J,..,IAR(10)}{Variable {\it I, J} are integers,
				   {\it IAR(10)} is array with 10 elements}
\fcc{\tab REAL A,X, ARRAY(-10:10)}{real numbers, array indices take values
$-10$, $-9$, $\ldots$, 9, 10}
\fcc{\tab DOUBLE PRECISION D,R(4,8)}{double precision,  
			{\tt R(4,8)} declares two dimensional
				array of 4 by 8}
\fcc{\tab COMPLEX  COMP}{complex number, the value is, e.g., (12.0, 0.7)}
\fcc{\tab LOGICAL LG}{logical variables take value {\tt .TRUE.} and
                                        {\tt .FALSE.}}
\fcc{\tab CHARACTER*5 CHAR}{character string of length 5}
\fcc{\tab PARAMETER ( PI = 3.14159 )}{declare constant}
\fcc{\tab DATA I/1/, A/1.2/, D/1.0D0/, CHAR/'WORDS'/, LG/.TRUE./}
			{set initial values}
\fcc{\tab COMMON/CNAME/X,R,...}{common blocks share variables between 
				subroutines}

\sect{Comparison}
\settabs 5 \columns
\+\fc{.EQ. (==)}{=} &
\fc{.LT. (<)}{<} &
\fc{.GT. (>)}{>} &
\fct{.AND.}{and} &
\fct{.OR.}{or} & \cr
\+\fc{.NE. (/=)}{\neq} &
\fc{.LE. (<=)}{\leq} &
\fc{.GE. (>=)}{\geq} &
\fct{.NOT.}{not} &
\fct{.XOR.}{exclusive or} & \cr
\+ & & & 
\fct{.EQV.}{equivalent} &
\fct{.NEQV.}{nonequivalent} & \cr

\sect{Flow Control}
\fcc{\tab DO 10 I = 1, 5}{do-loop with {\tt I} varies from 1 to 5 in step 1}
\fcc{10\ \ \ \ IAR(I) = I**2}{ }
\medskip
\fcc{\tab DO J = 0, 10, 2}{\logo do-loop with {\tt J} varies from 0 to 10 in step 2}
\fcc{\tab ...}{ }
\fcc{\tab END DO}{ }
\medskip
\medskip
%\vfill\break
\fcc{\tab DO}{\logo loop with exit condition specified anywhere in the block}
\shrink
\fcc{\tab \ \ ...}{ }
\fcc{\tab \ \ IF (...) EXIT}{loop terminate if condition is {\tt .TRUE.}}
\shrink
\fcc{\tab \ \ ...}{ }
\fcc{\tab END DO}{ }
\medskip
\medskip
\fcc{\tab DO WHILE ($expr$)}{\logo execute the body if $expr$ is {\tt .TRUE.}}
\fcc{\tab \ \ ...}{ }
\fcc{\tab END DO}{ }
\medskip
\fcc{\tab IF (J.GT.0) X = - 5}{conditional, assign {\tt X} to 5 if {\tt J > 0}}
\medskip
%%%%%%%%%%%%%%%%% FORCED PAGE BREAK %%%%%%%%%%%
\vfill\eject
%%%%%%%%%%%%%%%%%%%%%%%%%%%%%%%%%%%%%%%%%%%%%%%
\fcc{\tab IF ( ... ) THEN}{block if}
\shrink
\fcc{\tab \ \ ...}{ }
\fcc{\tab ELSE IF (...) THEN}{optional, several {\tt ELSE IF}s are allowed}
\shrink
\fcc{\tab \ \ ...}{ }
\fcc{\tab ELSE}{optional}
\shrink
\fcc{\tab \ \ ...}{ }
\fcc{\tab END IF}{ }
\medskip
\medskip
\fcc{\tab SELECT CASE ($expr$)}{\logo $expr$ may be integer, logical or
				character variable}
\fcc{\tab \ \ CASE ($value$)}{$value$ is constant expression}
\fcc{\tab \ \ \ \ $statement$}{}
\fcc{\tab \ \ CASE ($value_1$, $value_2$)}{}
\shrink
\fcc{\tab \ \ \ \ ...}{}
\fcc{\tab \ \ CASE DEFAULT}{}
\shrink
\fcc{\tab \ \ \ \ ...}{}
\fcc{\tab END SELECT}{}
\medskip
\fcc{\tab GOTO 15}{jump to line labelled as 15}

\sect{Input/Output}
\fcc{\tab OPEN(8,FILE='FILENAME',STATUS='UNKNOWN')}{associate unit
			number 8 with a file}
\fcc{\tab ...}{unit 5 is standard input (keyboard), 
				6 standard output (screen)}
\fcc{\tab READ(8,*) I, R, A}{ read from file, {\tt *} for free format}
\fcc{\tab WRITE(8,100) I, R, A}{ write to file}
\fcc{100\ \ \ FORMAT(1X,I5,F7.3,A1)}{write in the form of \]\]\]\]\],
			\]\]\]{\tt .}\]\]\], and \], respectively}
\fcc{\tab ...}{{\tt X} for space, {\tt I} for integer, {\tt F} for real,
		and {\tt A} for character}
\fcc{\tab CLOSE(8)}{close the opening file}

\sect{Character Strings}
\fcc{\tab 'a'//'bc'}{concatenate two strings, producing {\tt 'abc'}}
\fcc{\tab CS(2:5)}{extract substring in position 2 to 5}
\fcc{\tab LEN(CS)}{a function return the length of the string {\tt CS}}
\fcc{\tab ICHAR($char$), CHAR($integer$)}{convert character to integer,
			integer to character, respectively}

\sect{Intrinsic Functions}
\+\fcc{SIN(X)}{} &
\fcc{SINH(X)}{} &
\fcc{LOG(X)}{} &
\fcc{EXP(X)}{} &
\fcc{REAL(I)}{} & \cr
\+\fcc{COS(X)}{} &
\fcc{COSH(X)}{} &
\fcc{LOG10(X)}{} &
\fcc{INT(X)}{} &
\fcc{DFLOAT(I)}{} & \cr
\+\fcc{ABS(X)}{} &
\fcc{SQRT(X)}{} &
\fcc{MOD(I,J)}{} &
\fcc{TAN(X)}{}& 
\fcc{TANH(X)}{}& \cr
\+\fcc{NOT(I)}{} &
\fcc{IAND(I,J)}{} &
\fcc{IOR(I,J)}{} &
\fcc{IEOR(I,J)}{} &
\fcc{ISHFT(I,SHFT)}{} & \cr


\sect{Miscellaneous}
\fcc{\tab STOP 'Message'}{terminate execution with optional message}
\fcc{\tab RETURN}{return to calling program}
\fcc{\tab PAUSE}{pause}

\sect{New Fortran 90 features}
\medskip
\fcc{\tab TYPE INFO}{user defined type (structure) called {\tt INFO}}
\fcc{\tab \ \ CHARACTER*10 NAME}{ }
\fcc{\tab \ \ INTEGER AGE}{ }
\fcc{\tab END TYPE}{}
\medskip
\fcc{\tab TYPE (INFO) :: A, B, C}{define {\tt A,B,C} as type {\tt INFO}}
\fcc{\tab A\%NAME = 'js' }{refer the component by percent sign}
\medskip
\fcc{\tab REAL, TARGET :: D}{define {\tt D} so that address can be accessed by a pointer}
\fcc{\tab REAL, POINTER :: P}{pointer to a real number}
\fcc{\tab P => D}{pointer {\tt P} pointing to variable {\tt D}}
\fcc{\tab ALLOCATE(P); DEALLOCATE(P)}{dynamically allocate/free memory}
\medskip
\fcc{\tab MODULE $name$}{module is a piece of variable declarations and codes}
\fcc{\tab \ \ USE Others}{which can be used by other module or subprograms}
\fcc{\tab \ \ {\it variable declarations}}{global variables inside module}
\fcc{\tab \ \ CONTAINS}{subroutine or function definition follows}
\fcc{\tab \ \ {\it subprograms} }{}
\fcc{\tab END MODULE}{}
\medskip

\vfill
\eject
\end
